\section{Branched Covers and the Genus-Zero Classification}
\label{sec:branched-covers-genus-zero}

\begin{layman}
A nonconstant holomorphic map between compact Riemann surfaces is
never just an injection --- it is a \emph{branched covering}. Most
points in the target have the same finite number \(d\) of preimages
in the source; at finitely many \emph{branch points} the preimages
collide with multiplicity. The number \(d\) is the \emph{degree} of
the map, the basic numerical invariant of the cover. A degree-one
holomorphic map between compact Riemann surfaces --- generic fibre is
a single point --- is automatically a homeomorphism, and in fact a
biholomorphism: this is one of the chapter's punch lines and the
key step in the genus-zero classification.

The chapter begins by formalising what ``degree'' means. Each
preimage \(p\) of a base point \(q\) carries a \emph{ramification
index} \(e_p(f)\), the local-power-series order of \(f - f(p)\) at
\(p\). The fibre sum \(\sum_{p \in f^{-1}(q)} e_p(f)\) is the
weighted count of preimages; the weighted-fibre-sum-is-locally-
constant theorem then forces this count to be the same for every
\(q\) --- that constant value is the degree. Showing local constancy
costs real work: the unramified locus is open, ramified points are
isolated and finite, the chart-local power series can be put in a
local conjugate-to-\(z^k\) normal form via a holomorphic \(k\)-th
root, and the weighted sum is locally constant on each side and
glues across.

With the degree definition in hand, the chapter develops the
infrastructure that pins down a degree-one map. At every unramified
point the map is locally a chart-to-chart inverse; the entire fibre
of a regular value is a singleton when degree is one; an unramified
nonconstant holomorphic map is a local biholomorphism; a degree-one
nonconstant holomorphic map between compact Riemann surfaces is a
homeomorphism. Combine with Riemann-Roch on a genus-zero surface
(producing a degree-one meromorphic function to the Riemann sphere)
and you get the first anti-hack theorem of the challenge:
\emph{analytic genus zero is equivalent to being homeomorphic to
the 2-sphere}. The chapter ends with that classification.

Why bother? The genus-zero classification is one of the four anti-
hack theorems pinned by \texttt{Jacobian/Challenge.lean}. It rules
out trivial fake constructions like \texttt{genus X := 0}: the
challenge demands that \texttt{genus X = 0} actually entails \(X\)
is the topological sphere. Reaching that conclusion goes through
exactly this chapter --- branched degree, Riemann-Hurwitz, degree-one
classification, and the analytic-genus-zero \(\Rightarrow\)
sphere argument.
\end{layman}


\begin{definition}[Branched-cover degree of a nonconstant holomorphic map]
\label{def:branched-degree}
\uses{def:vanishing-order}
For a nonconstant holomorphic \(f:X\to Y\), the degree \(\deg(f)\) is
the cardinality, counted with multiplicity, of any generic fiber. Each
preimage \(p\in f^{-1}(q)\) contributes its ramification index
\(e_p(f)=\ord_p(f - f(p))\ge 1\); the sum is constant in \(q\) and
equals \(\deg(f)\).
\depends{Local power-series order at \(p\) (\texttt{analyticOrderAt});
constancy in \(q\) is the content of
\texttt{thm:branched-degree-well-defined}, decomposed into
\texttt{lem:branch-locus-finite},
\texttt{lem:local-inj-unramified},
\texttt{lem:local-kfold-ramified}, and
\texttt{lem:weighted-fiber-locally-constant}.}
\lean{JacobianChallenge.HolomorphicForms.branchedDegree}
\leanok
\end{definition}

\begin{layman}
For a nonconstant holomorphic map between two compact Riemann surfaces,
the \emph{degree} counts how many times the map covers its target.
Pick a generic point downstairs and count its preimages, with
multiplicities at branch points where two or more sheets pinch together:
that count is the degree, and it's the same whichever generic target
point you choose. The multiplicity at a preimage is its
\emph{ramification index}, which measures how many sheets meet there. An
index of \(1\) means the map is a clean local homeomorphism nearby; an
index of \(2\) means two sheets pinch into one (like the square-root
function); higher indices mean more pinching. Why bother? Degree turns
holomorphic maps between surfaces into measurable things --- the way
``2-to-1 covering'' or ``5-to-1 covering'' captures the structure. Most
of the classical theory of maps between Riemann surfaces (Riemann--Hurwitz,
the trace, push-pull) runs on this number.
\end{layman}


\begin{proof}\leanok\end{proof}


\begin{proof}\leanok\end{proof}


\begin{proof}\leanok\end{proof}

\begin{layman}
The branch points of a map --- places where two or more sheets pinch
together, so the ramification index jumps above \(1\) --- can never form
an infinite cluster. Locally, ramification means the derivative of the
map vanishes. Derivatives of holomorphic gadgets are themselves
holomorphic, and on a compact surface the zero set of a nonzero
holomorphic gadget is always discrete and finite, by the same
identity-theorem-plus-compactness combo we saw for divisors. So branch
points come in tidy finite quantity. Why bother? Most theorems about
branched covers are easier to state and prove ``away from the branch
locus,'' and knowing that locus is finite guarantees we can always
extend conclusions across it by continuity --- exactly how the trace of a
1-form will be defined later.
\end{layman}

\begin{proofsketch}
\(e_p(f)\ge 2\) iff the local derivative of \(f\) vanishes. The
derivative is itself meromorphic; its zero set is discrete by the
identity theorem, and finite by compactness.
\end{proofsketch}

\begin{lemma}[Local injectivity at an unramified point]
\label{lem:local-inj-unramified}
\uses{def:branched-degree, lem:order-eq-one-iff-deriv-ne-zero, lem:chart-local-inverse-unramified}
If \(e_x(f)=1\) at a point \(x\in X\), there is an open neighborhood
\(U\ni x\) and an open neighborhood \(V\ni f(x)\) such that for every
\(y\in V\), the set \(U\cap f^{-1}(y)\) is a singleton.
\depends{Sub-sub-leaves: \texttt{lem:order-eq-one-iff-deriv-ne-zero}
(common B1=A1 helper) and \texttt{lem:chart-local-inverse-unramified}
(B2, Mathlib's \texttt{AnalyticAt.localInverse}); assembly is chart
transport.}
\lean{JacobianChallenge.Blueprint.IsHolomorphicAt.exists_local_inj_of_unramified}
\leanok
\end{lemma}
\begin{proof}\leanok\end{proof}

\begin{lemma}[Order one iff chart-local derivative nonzero]
\label{lem:order-eq-one-iff-deriv-ne-zero}
\uses{def:branched-degree, def:vanishing-order}
The chart-local order \(e_x(f)=1\) iff the derivative
\((\chi\circ f\circ\varphi^{-1})'(\varphi(x))\ne 0\), where
\((\varphi,\chi)\) are the canonical charts at \(x\) and \(f(x)\).
\depends{Mathlib's
\texttt{AnalyticAt.analyticOrderAt\_sub\_eq\_one\_of\_deriv\_ne\_zero}
plus the
\texttt{analyticOrderNatAt}\(\,=1\) \(\iff\) \texttt{analyticOrderAt}\(\,=1\) characterisation.}
\lean{JacobianChallenge.Blueprint.mapAnalyticOrderAt_eq_one_iff_chartLocal_deriv_ne_zero}
\leanok
\end{lemma}
\begin{proof}\leanok\end{proof}

\begin{lemma}[Chart-local inverse at unramified]
\label{lem:chart-local-inverse-unramified}
\uses{lem:order-eq-one-iff-deriv-ne-zero}
At an unramified point \(x\), the chart-local presentation
\(\chi\circ f\circ\varphi^{-1}\) admits a holomorphic local inverse
near \(\chi(f(x))\).
\depends{\texttt{AnalyticAt.localInverse} from
\texttt{Mathlib/Analysis/Analytic/Inverse.lean}.}
\lean{JacobianChallenge.Blueprint.chartLocalAt_localInverse_of_unramified}
\leanok
\end{lemma}
\begin{proof}\leanok\end{proof}

\begin{layman}
At a point where the holomorphic map doesn't pinch sheets together
(ramification index \(1\)), the map is a clean local
homeomorphism: zoom in enough and you see a single sheet covering a
single small disk, with each downstairs point having a unique upstairs
preimage. Why bother? This is the building block for conservation of
the weighted fibre count off branch values: every unramified preimage
contributes exactly \(1\) to the count, and that count is locally
stable because nearby downstairs points have exactly the same
unique-preimage neighborhoods.
\end{layman}

\begin{proofsketch}
The chart-local power-series expansion at \(x\) is
\(f(t)=a\cdot t+O(t^2)\) with \(a\ne 0\), so the chart-local derivative
is nonzero at the chart image of \(x\). Apply the analytic inverse
function theorem for \(\C\to\C\) to get a chart-local inverse on a
neighborhood; transport back through the canonical charts on \(X\)
and \(Y\).
\end{proofsketch}

\begin{lemma}[Local power-series form at order \texorpdfstring{\(k\)}{k}]
\label{lem:chart-local-pow-mul}
\uses{def:branched-degree, def:vanishing-order}
If \(e_x(f)=k\ge 1\), the centred chart-local difference factors as
\((t-z_0)^k\cdot g(t)\) for some analytic \(g\) with \(g(z_0)\ne 0\),
where \(z_0\) is the chart image of \(x\).
\depends{Direct application of Mathlib's
\texttt{AnalyticAt.analyticOrderAt\_eq\_natCast}.}
\lean{JacobianChallenge.Blueprint.chartLocalAt_eq_pow_mul_of_order}
\leanok
\end{lemma}
\begin{proof}\leanok\end{proof}

\begin{lemma}[Holomorphic \texorpdfstring{\(k\)}{k}-th root]
\label{lem:analytic-kth-root}

If \(g\) is analytic at \(z_0\) with \(g(z_0)\ne 0\) and \(k\ge 1\),
there is an analytic \(h\) near \(z_0\) with \(h(z)^k=g(z)\) locally.
\depends{Holomorphic logarithm
(\texttt{Complex.analyticAt\_log} on a simply-connected nbhd of
\(g(z_0)\) avoiding \(0\)) plus \(\exp\) plus
\(h:=\exp((1/k)\cdot\log g)\).}
\lean{JacobianChallenge.Blueprint.analyticAt_kth_root_of_ne_zero}
\leanok
\end{lemma}
\begin{proof}\leanok\end{proof}

\begin{lemma}[Local \texorpdfstring{\(k\)}{k}-fold structure at a ramified point]
\label{lem:local-kfold-ramified}
\uses{def:branched-degree, def:vanishing-order, lem:chart-local-pow-mul, lem:analytic-kth-root}
If \(e_x(f)=k\ge 1\), there is an open neighborhood \(U\ni x\) and an
open neighborhood \(V\ni f(x)\) such that for every \(y\in V\) with
\(y\ne f(x)\), the set \(U\cap f^{-1}(y)\) consists of exactly \(k\)
distinct unramified preimages.
\depends{Sub-sub-leaves: \texttt{lem:chart-local-pow-mul} (C1, local
power-series form) and \texttt{lem:analytic-kth-root} (C2,
holomorphic \(k\)-th root); assembly is the substitution
\(s=(t-z_0)\cdot h(t)\) flattening the local form to \(s\mapsto s^k\),
which has exactly \(k\) simple preimages of any nonzero target.}
\lean{JacobianChallenge.Blueprint.IsHolomorphicAt.exists_local_kfold_of_ramified}
\leanok
\end{lemma}
\begin{proof}\leanok\end{proof}

\begin{layman}
At a branch point of order \(k\), the map locally looks like
\(z\mapsto z^k\): the \(k\) sheets that meet at the branch point pull
apart as you move slightly away, becoming \(k\) distinct simple
preimages of any nearby target. This is what conserves the weighted
fibre count across branch values: the single ramified preimage with
weight \(k\) splits into \(k\) simple preimages, each with weight
\(1\), summing to \(k\) again.
\end{layman}

\begin{proofsketch}
Write \(f(t)=t^k\cdot g(t)\) in a chart at \(x\), with \(g(0)\ne 0\).
Choose a holomorphic \(k\)-th root \(g(t)^{1/k}\) (exists since \(g\)
is nonvanishing near \(0\)) and substitute \(s=t\cdot g(t)^{1/k}\) to
flatten the local form to \(f(s)=s^k\). The map \(s\mapsto s^k\) has
exactly \(k\) simple preimages of any nonzero point. Transport back
through charts.
\end{proofsketch}

\begin{lemma}[Properness of fibre below an open neighborhood]
\label{lem:proper-fiber-shrink}

For a continuous map \(f\) from a compact source to a Hausdorff
target, and any open \(U\) containing \(f^{-1}(y_0)\), there is a
neighborhood \(V\) of \(y_0\) such that for every \(y\in V\),
\(f^{-1}(y)\subseteq U\).
\depends{\(U^c\) is closed-in-compact \(\Rightarrow\) compact; image \(f(U^c)\) is
compact \(\Rightarrow\) closed in \(T_2\); \(y_0\notin f(U^c)\); take
\(V:=f(U^c)^c\).}
\lean{JacobianChallenge.Blueprint.eventually_fiber_subset_of_compact_T2}
\leanok
\end{lemma}
\begin{proof}\leanok\end{proof}

\begin{lemma}[Weighted fibre count locally constant at a single point]
\label{lem:weighted-fiber-eventually-eq}
\uses{lem:local-inj-unramified, lem:local-kfold-ramified, lem:proper-fiber-shrink}
At every \(y_0\in Y\), there is a neighborhood \(V\ni y_0\) on which
the weighted fibre sum is constant equal to its value at \(y_0\).
\depends{Choose pairwise-disjoint open nbhds \(U_x\) around each
\(x\in f^{-1}(y_0)\) (T2 + finite); shrink each \(U_x\) so that
either \texttt{lem:local-inj-unramified} (if \(e_x=1\)) or
\texttt{lem:local-kfold-ramified} (if \(e_x\ge 2\)) applies, with an
associated \(V_x\); intersect with the neighborhood from
\texttt{lem:proper-fiber-shrink} to ensure all preimages of nearby
\(y\) lie in \(\bigcup U_x\); sum contributions.}
\lean{JacobianChallenge.Blueprint.weightedFiberSum_eventually_eq}
\leanok
\end{lemma}
\begin{proof}\leanok\end{proof}

\begin{lemma}[Weighted fibre count is locally constant]
\label{lem:weighted-fiber-locally-constant}
\uses{lem:weighted-fiber-eventually-eq}
For a nonconstant holomorphic \(f:X\to Y\) between a compact
preconnected complex \(1\)-manifold and a Hausdorff complex
\(1\)-manifold, the weighted fibre sum
\(y\mapsto\sum_{p\in f^{-1}(y)}e_p(f)\) is locally constant on \(Y\).
\depends{Direct corollary of
\texttt{lem:weighted-fiber-eventually-eq} (which holds at every
\(y_0\)) plus
\texttt{IsLocallyConstant.iff\_eventuallyEq} from Mathlib.}
\lean{JacobianChallenge.Blueprint.isHolomorphic_weightedFiberSum_isLocallyConstant}
\leanok
\end{lemma}
\begin{proof}\leanok\end{proof}

\begin{layman}
The weighted fibre count is the same for nearby downstairs points: at
each preimage, the local picture (either a single unramified sheet of
weight \(1\), or \(k\) simple sheets summing to \(k\), or one
ramified sheet of weight \(k\)) contributes the same amount whether
you're sitting exactly at the branch value or slightly displaced.
Stitching together all the preimages of a target point --- there are
finitely many --- and choosing disjoint neighborhoods around each, the
total weighted count is locally constant.
\end{layman}

\begin{theorem}[Branched degree is well-defined]
\label{thm:branched-degree-well-defined}
\uses{lem:weighted-fiber-locally-constant}
For a nonconstant holomorphic \(f:X\to Y\) between compact preconnected
complex \(1\)-manifolds, the weighted fibre sum is constant on \(Y\):
the same nonnegative integer \(\deg(f)\) is obtained at every base
point.
\depends{Local-constancy from
\texttt{lem:weighted-fiber-locally-constant} plus
\texttt{IsLocallyConstant.apply\_eq\_of\_preconnectedSpace} from
Mathlib (preconnected target \(\Rightarrow\) globally constant).}
\lean{JacobianChallenge.Blueprint.isHolomorphic_weightedFiberSum_const}
\leanok
\end{theorem}
\begin{proof}\leanok\end{proof}

\begin{layman}
Locally constant on a connected target means globally constant. So the
weighted fibre count assigns the same number --- the degree --- to every
target point. This is the proof that ``the degree of a holomorphic map''
is actually a single well-defined number, not a value that drifts as
you move around the target.
\end{layman}



\begin{layman}
If a holomorphic map between compact Riemann surfaces has degree one,
every fiber's ramification indices sum to one. Since each index is at
least one, the only way the total can equal one is if the fiber consists
of exactly one point with index exactly one. So no branching anywhere ---
the map is everywhere a clean local homeomorphism. Why bother?
Degree-one maps are the candidate isomorphisms between Riemann surfaces.
Showing they have no branching is the first hurdle; the next hurdle
(local biholomorphism) and the final hurdle (global bijection) follow
from there. The whole genus-zero classification runs on this chain.
\end{layman}

\begin{proof}\leanok
At a generic fiber, the sum of ramification indices equals \(1\), so the
fiber has exactly one point with index \(1\). At a branch point \(p\),
the sum of indices over the fiber is also \(1\), but each index is
\(\ge 1\), so the fiber is a singleton \(\{p\}\) with \(e_p(f)=1\).
\end{proof}


\begin{proof}\leanok\end{proof}

\begin{layman}
Take a nonconstant holomorphic map that's unramified --- meaning it never
pinches sheets together. Around every point, the map is a \emph{local
biholomorphism}: a small neighborhood of the source maps in a smooth,
invertible, holomorphic way onto a small neighborhood of the target.
The mechanism is the inverse function theorem from one-variable complex
analysis: unramified means the local derivative is nonzero, and a
nonzero derivative is exactly the condition for a holomorphic function
to have a holomorphic local inverse. Why bother? Locally, the map looks
like the identity (after a change of coordinates). Combined with
degree-one (every fiber a singleton), this will force the whole map to
be a global biholomorphism --- the isomorphism we need to identify
genus-zero surfaces with the sphere.
\end{layman}

\begin{theorem}[Degree-one is bijective]
\label{thm:degree-one-bijective}
\uses{def:branched-degree}
A degree-one holomorphic map \(f:X\to Y\) is bijective.
\lean{JacobianChallenge.HolomorphicForms.degree_one_bijective}
\leanok
\end{theorem}
\begin{proof}\leanok\end{proof}

\begin{layman}
A degree-one holomorphic map between compact Riemann surfaces is a
bijection. \emph{Surjective}: a nonconstant holomorphic map between
compact surfaces is always onto --- its image is open (open mapping
theorem) and closed (compactness), and a nonempty subset that's both
open and closed in a connected space is the whole thing. \emph{Injective}:
each fiber has size exactly one, by the no-ramification lemma. So the
map is both onto and one-to-one --- a set-theoretic bijection. Why bother?
A continuous bijection between compact Hausdorff spaces is automatically
a homeomorphism (next theorem), and combined with the local
biholomorphism from above, that promotes our degree-one map to a full
biholomorphism between surfaces.
\end{layman}

\begin{proofsketch}
Surjective: a degree-\(d\) holomorphic map between compact Riemann
surfaces is surjective (the image is open by open-mapping and closed by
compactness). Injective: by
Lemma~\ref{lem:degree-one-no-ramification}, every fiber has exactly one
point with ramification index 1, hence cardinality 1.
\end{proofsketch}


\begin{proof}\leanok\end{proof}

\begin{layman}
A purely topological gem: any continuous one-to-one onto map from one
compact Hausdorff space to another is automatically a homeomorphism ---
its inverse is continuous for free, no extra checking required. The
reason is a chain of clean implications. Closed subsets of a compact
space are themselves compact. Compact subsets of a Hausdorff space are
closed. So the map sends closed sets to closed sets, which is exactly
continuity of the inverse. Why bother? Without it we'd have to manually
verify that the inverse is continuous everywhere, which on a curvy
surface is fiddly. With it, going from ``continuous bijection'' to
``isomorphism of topological spaces'' is one line. Mathlib already has
this packaged.
\end{layman}


\begin{proof}\leanok\end{proof}

\begin{layman}
We package the chain ``degree-one nonconstant holomorphic map between
compact Riemann surfaces is a biholomorphism'' as a single classical
input. It bundles three steps: degree one means no ramification; no
ramification means a local biholomorphism everywhere; and a continuous
bijection between compact Hausdorff spaces is a homeomorphism --- together
giving a global holomorphic isomorphism. Why bother? Combined with the
genus-zero existence proposition (genus zero gives a degree-one map to
the sphere), this umbrella delivers the forward direction of the
genus-zero classification: any genus-zero compact Riemann surface is
\emph{biholomorphic} to the Riemann sphere, hence in particular
homeomorphic to the topological sphere. That's the anti-hack
``genus zero iff homeomorphic to sphere'' baked into the challenge
statement.
\end{layman}

\begin{proof}[Proof of Theorem~\ref{thm:genus-zero-classification}]
If \(g(X)=0\), Proposition~\ref{prop:genus-zero-degree-one-map} and
Input~\ref{input:degree-one-isomorphism} identify \(X\) biholomorphically,
hence homeomorphically, with \(\CPone\), whose underlying topological space is
\(\Sphere\).

Conversely, if \(X\) is homeomorphic to \(\Sphere\), the topological genus is
zero. The comparison theorem between analytic genus
\(\dim_{\C}H^{0}(X,\Omega^{1}_{X})\) and topological genus gives \(g(X)=0\).
\end{proof}

\begin{formalizationnote}
The current Lean target splits this into the two directions in
\code{Jacobian/HolomorphicForms/GenusZeroClassification.lean}.

\emph{Forward direction — analytic uniformization route (not Riemann--Roch).}
Although the classical argument above routes through Riemann--Roch
(genus \(0\Rightarrow\) a degree-one meromorphic map to \(\CPone\)), the Lean
development \emph{cannot} take that route: the Riemann--Roch section chain
(\code{genusZero\_pointRRSection\_*} in \code{RiemannRoch.lean} /
\code{MeromorphicToBranchedCover.lean}) is circular here --- its
\code{\_with\_analyticData} exit feeds back through the very genus-zero
single-pole data it is meant to produce. So the forward direction is instead
built from a self-contained \emph{local analytic uniformization substrate}:
\begin{enumerate}
\item local chart-ball power series for holomorphic chart maps
  (\code{DiffContOnCl.hasFPowerSeriesOnBall});
\item open-mapping / inverse-function local charts (nonzero derivative
  \(\Rightarrow\) locally injective, locally open, normalizable);
\item a normal-family (Montel / Arzel\`a--Ascoli) extraction selecting
  compatible normalized local coordinate patches;
\item gluing the patches through the two-chart \(\CPone\) atlas into a global
  homeomorphism \(X\simeq\CPone\) that is \code{ContMDiff} in both directions.
\end{enumerate}
This is packaged as the single high-level provider
\code{exists\_biholomorph\_onePoint\_of\_genus\_zero}
(\code{GenusZeroClassification.lean}), with
\code{exists\_biholomorph\_onePoint\_of\_analyticGenus\_zero} and
\code{complexSimplePolePrincipalPart\_of\_biholomorph\_onePoint}
(\code{MeromorphicToBranchedCover.lean}) the thin layers that recover the
analytic-genus hypothesis and pull back the standard simple-pole coordinate.
The simple-pole / degree-one meromorphic map is thus \emph{obtained from} the
biholomorphism, rather than used to construct it.

\emph{Reverse direction.} A surface homeomorphic to \(\Sphere\) has analytic
genus zero, via a Liouville-style computation that
\(H^{0}(\CPone,\Omega^{1})=0\) (a holomorphic 1-form is \(f(z)\,dz\) on the
identity chart; the inversion-chart transition \(w=z^{-1}\) sending it to
\(-f(1/w)/w^{2}\,dw\) forces \(f\equiv0\) by Liouville) transported across the
homeomorphism by
\code{analyticGenus\_eq\_zero\_of\_homeomorphic\_sphere}, plus comparison of
analytic and topological genus.
\end{formalizationnote}

\subsection{Identification with \(\CPone\): coordinate chart details}

\begin{lemma}[Continuity of local representatives]
\label{lem:section-localRepr-continuity}
For a holomorphic section of the cotangent bundle, the local representative in the identity and inversion charts are continuous.
\lean{JacobianChallenge.HolomorphicForms.ContMDiffSection_localRepr_identityChart_contDiff, JacobianChallenge.HolomorphicForms.ContMDiffSection_localRepr_inversionChart_continuousAt_zero}
\leanok
\end{lemma}

\begin{theorem}[Isomorphism with \(\CPone\) one-forms]
\label{lem:holomorphic-one-form-linearEquiv-to-OnePointCx}
\uses{lem:section-localRepr-continuity, lem:onepoint-cx-chart-overlap-pullback}
For any Riemann surface \(X\) biholomorphic to \(\CPone\), there is a linear equivalence between their spaces of holomorphic 1-forms.
\lean{JacobianChallenge.HolomorphicForms.holomorphicOneForm_linearEquiv_of_biholo_to_OnePointCx}
\leanok
\end{theorem}

\subsection{Degree-one meromorphic-to-sphere classification: structural decomposition}
\label{sec:degree-one-mer-decomp}

The chain ``meromorphic map with pole divisor \([P]\) \(\Rightarrow\)
continuous \(\Rightarrow\) bijective'' factors into structural axioms
on the project's abstract \code{MeromorphicMapToSphere} interface,
since this interface does not yet expose the geometric content
linking the abstract \code{poleDivisor} to the concrete \code{toMap}.

\begin{lemma}[Continuity off the pole]
\label{lem:meromorphic-continuous-off-pole}
\uses{lem:meromorphic-no-infty-off-pole, lem:meromorphic-continuous-on-finite-set}
A meromorphic map \(f:X\to\CPone\) with pole divisor \([P]\) is
continuous on \(X\setminus\{P\}\).
\lean{JacobianChallenge.HolomorphicForms.MeromorphicMapToSphere.continuousOn_compl_pole_of_poleDivisor_point}
\leanok
\end{lemma}
\begin{proof}\leanok
Off the pole locus, \(f\) takes only finite values. In any chart at
a finite point, \(f\) coincides with a smooth \(\C\)-valued
function (the chart-local representative). Smoothness implies
continuity in the chart, and the inclusion \(\iota:\C\to\CPone\) is
continuous. Hence \(f\) is continuous at every point of
\(X\setminus\{P\}\).
\end{proof}

\begin{lemma}[Tendency to \(\infty\) at a simple pole]
\label{lem:meromorphic-tends-to-infty-at-simple-pole}
\uses{lem:simple-pole-modulus-diverges, lem:onepoint-tendsto-infty-of-modulus-diverges, lem:meromorphic-value-at-pole}
A meromorphic map \(f:X\to\CPone\) with pole divisor \([P]\) tends
to \(\infty\in\CPone\) along the neighbourhood filter at \(P\).
\lean{JacobianChallenge.HolomorphicForms.MeromorphicMapToSphere.tendsto_infty_at_pole_of_poleDivisor_point}
\leanok
\end{lemma}
\begin{proof}\leanok
In a chart at \(P\), the local representative of \(f\) is a
meromorphic function with a simple pole at the chart image of \(P\).
The local Laurent expansion takes the form
\(z\mapsto a/(z-z_{0})+\text{holomorphic}\) with \(a\neq 0\), which
diverges in modulus as \(z\to z_{0}\). The one-point topology on
\(\CPone\) translates modulus-divergence into convergence to
\(\infty\).
\end{proof}

\begin{lemma}[Value at a simple pole is \(\infty\)]
\label{lem:meromorphic-value-at-pole}

A meromorphic map \(f:X\to\CPone\) with pole divisor \([P]\) has
\(f(P)=\infty\).
\lean{JacobianChallenge.HolomorphicForms.MeromorphicMapToSphere.toMap_pole_eq_infty_of_poleDivisor_point}
\leanok
\end{lemma}
\begin{proof}\leanok
The data of a meromorphic map records both the divisor and the
extended map. By construction, the extension assigns the value
\(\infty\) at every pole point. (At the project's interface level,
this is an axiomatic structural field that the eventual concrete
construction will discharge.)
\end{proof}

\begin{theorem}[Continuous extension over a simple pole]
\label{thm:meromorphic-continuous-extension}
\uses{lem:meromorphic-continuous-off-pole, lem:meromorphic-tends-to-infty-at-simple-pole, lem:meromorphic-value-at-pole}
A meromorphic map \(f:X\to\CPone\) with pole divisor \([P]\) on a
compact Riemann surface is continuous as a map \(X\to\CPone\).
\lean{JacobianChallenge.HolomorphicForms.meromorphicMapToSphere_continuous_of_poleDivisor_point}
\leanok
\end{theorem}
\begin{proof}\leanok
Continuity is local. At any point \(x\neq P\) use
Lemma~\ref{lem:meromorphic-continuous-off-pole}. At \(x=P\), the
limit characterisation of continuity reduces to
\(\lim_{x\to P}f(x)=f(P)\), which is
Lemma~\ref{lem:meromorphic-tends-to-infty-at-simple-pole} together
with Lemma~\ref{lem:meromorphic-value-at-pole}.
\end{proof}

\begin{lemma}[Bridge: meromorphic-to-branched-cover-data]
\label{lem:meromorphic-to-branched-cover-data}
\uses{def:branched-degree, thm:meromorphic-continuous-extension}
A meromorphic map \(f:X\to\CPone\) with \(\deg(\operatorname{poles}(f))=1\) admits a
\code{BranchedCoverData} of branched degree 1.
\lean{JacobianChallenge.HolomorphicForms.MeromorphicMapToSphere.exists_branchedCoverData_of_pole_degree_one}
\leanok
\end{lemma}
\begin{proof}
\leanok
Extract the structural axiom field
\code{hasBranchedCoverDataOfPoleDegree} from the
\code{MeromorphicMapToSphere} carrier, which (conditional on
\code{Continuous toMap}) supplies a \code{BranchedCoverData} with
\code{branchedDegree = poleDivisor.degree.toNat}; substitute the
pole-degree-1 hypothesis to land at \code{branchedDegree = 1}.
\end{proof}

% NOTE (audit): Earlier iterations introduced
% lem:degree-one-surjective + lem:degree-one-injective at the
% MeromorphicMapToSphere level. After cross-section audit, the existing
% Sec02 thm:degree-one-bijective (already \leanok) covers the same
% mathematical content at the BranchedCoverData level. The
% MeromorphicDegree.lean Lean is now refactored to route M4/M5 through
% the bridge `lem:meromorphic-to-branched-cover-data` above into
% `Blueprint.degree_one_bijective`, eliminating the parallel chain.

\begin{theorem}[Degree-one bijectivity (meromorphic level)]
\label{thm:degree-one-bijectivity-meromorphic}
\uses{lem:meromorphic-to-branched-cover-data, thm:degree-one-bijective}
A continuous meromorphic map \(f:X\to\CPone\) of pole-degree one is
bijective.
\lean{JacobianChallenge.HolomorphicForms.meromorphicMapToSphere_bijective_of_poleDivisor_degree_one}
\leanok
\end{theorem}
\begin{proof}\leanok
Bridge through Lemma~\ref{lem:meromorphic-to-branched-cover-data} to a
\code{BranchedCoverData} of branched degree 1, then apply
Theorem~\ref{thm:degree-one-bijective} (Sec02, already \(\leanok\)).
\end{proof}

\subsection{Uniformization-lite: holomorphic 1-forms vanish on the topological sphere}
\label{sec:uniformization-lite}

\begin{lemma}[Holomorphic 1-form equivalence under sphere homeomorphism]
\label{lem:holomorphic-one-form-equiv-of-homeo-sphere}
\uses{thm:uniformization-genus-zero-biholomorphism, lem:holomorphic-one-form-pullback-via-biholo}
If a compact connected complex 1-manifold \(X\) is homeomorphic
(as a topological space) to the standard \(2\)-sphere, then there
is a \(\C\)-linear equivalence
\(H^{0}(X,\Omega^{1}_{X})\cong H^{0}(\CPone,\Omega^{1}_{\CPone})\).
\lean{JacobianChallenge.HolomorphicForms.holomorphicOneForm_linearEquiv_of_homeoSphere_exists}
This transport step depends on the open genus-zero uniformization frontier:
Theorem~\ref{thm:uniformization-genus-zero-biholomorphism} upgrades the
sphere homeomorphism to a biholomorphism \(X\simeq\code{OnePoint}\,\C\).
Only after that smooth/biholomorphic upgrade does
Lemma~\ref{lem:holomorphic-one-form-pullback-via-biholo} supply the pullback
linear equivalence of holomorphic \(1\)-form spaces.
\end{lemma}

\begin{theorem}[Holomorphic 1-forms vanish on the topological sphere]
\label{thm:homeo-sphere-holomorphic-one-form-vanishing}
\uses{lem:holomorphic-one-form-equiv-of-homeo-sphere, lem:onepoint-holomorphic-one-form-subsingleton}
A compact connected Riemann surface homeomorphic to the standard
\(2\)-sphere has \(\dim_{\C}H^{0}(X,\Omega^{1}_{X})=0\); equivalently,
the space of holomorphic 1-forms is a subsingleton.
\lean{JacobianChallenge.HolomorphicForms.homeoSphereHolomorphicOneFormVanishing}
By Lemma~\ref{lem:holomorphic-one-form-equiv-of-homeo-sphere}, the
holomorphic 1-form space on \(X\) is \(\C\)-linearly isomorphic to
that on \(\CPone\). Lemma~\ref{lem:onepoint-holomorphic-one-form-subsingleton}
shows that the latter is a subsingleton. Subsingleton transports through
linear equivalences, so \(H^{0}(X,\Omega^{1}_{X})\) is itself a subsingleton.
\end{theorem}

\subsection{Identity- and inversion-chart coefficient extraction}
\label{sec:chart-coeff-extraction}

\begin{lemma}[Identity-chart coefficient is \(C^{\infty}\)]
\label{lem:identity-chart-coeff-contdiff}
\uses{lem:section-localRepr-identity-chart-contdiff}
For \(\omega\in H^{0}(\CPone,\Omega^{1}_{\CPone})\), the function
\(\C\to\C\) obtained by reading \(\omega\) in the identity chart
of \(\CPone\) and evaluating on \(1\in\C\) is \(C^{\infty}\).
\lean{JacobianChallenge.HolomorphicForms.holomorphicOneFormIdentityChartCoeffContDiff}
\leanok
\end{lemma}
\begin{proof}\leanok
The cotangent-bundle section \(\omega\) is smooth in any chart of
\(\CPone\). In the identity chart, the local representative is a
smooth section of the trivial cotangent bundle on \(\C\), i.e.\ a
smooth covector field. Evaluating a smooth covector field on a
constant tangent vector yields a smooth scalar function. The
specific Mathlib gap is the \emph{chart-trivialisation API for
\code{ContMDiffSection}} on the cotangent bundle, which is absent
in v4.28.0 --- once supplied, the proof is immediate.
\end{proof}

\begin{lemma}[Direct finite-point coefficient is \(C^{\infty}\)]
\label{lem:holomorphic-one-form-coeff-contdiff}
\uses{lem:identity-chart-coeff-contdiff}
For \(\omega\in H^{0}(\CPone,\Omega^{1}_{\CPone})\), the direct finite-point
coefficient \(z\mapsto \omega_{z}(1)\) is \(C^{\infty}\) as a function
\(\C\to\C\).
\lean{JacobianChallenge.HolomorphicForms.holomorphicOneFormCoeffContDiff}
\leanok
\end{lemma}
\begin{proof}\leanok
Lemma~\ref{lem:identity-chart-coeff-contdiff} gives smoothness of the
coefficient read through the identity-chart local representative. The identity
chart on finite points has inverse \(z\mapsto z\in\code{OnePoint}\,\C\), so this
chart-local coefficient agrees with the direct coefficient used in the
Liouville assembly.
\end{proof}

\begin{lemma}[Finite-point coefficient is entire]
\label{lem:holomorphic-one-form-coeff-entire}
\uses{lem:holomorphic-one-form-coeff-contdiff}
For \(\omega\in H^{0}(\CPone,\Omega^{1}_{\CPone})\), the direct finite-point
coefficient \(z\mapsto \omega_{z}(1)\) is complex differentiable on all of
\(\C\).
\lean{JacobianChallenge.HolomorphicForms.holomorphicOneForm_coeff_entire}
\leanok
\end{lemma}
\begin{proof}\leanok
This is the standard consequence of
Lemma~\ref{lem:holomorphic-one-form-coeff-contdiff}: a \(C^{\infty}\) complex
function is complex differentiable everywhere.
\end{proof}

\begin{lemma}[Inversion-chart coefficient is continuous at \(0\)]
\label{lem:inversion-chart-coeff-continuous-at-zero}
\uses{lem:section-localRepr-inversion-chart-continuous-at-zero}
For \(\omega\in H^{0}(\CPone,\Omega^{1}_{\CPone})\), the function
\(\C\to\C\) obtained by reading \(\omega\) in the inversion chart
of \(\CPone\) and evaluating on \(1\in\C\) is continuous at
\(0\in\C\).
\lean{JacobianChallenge.HolomorphicForms.holomorphicOneFormInversionChartCoeffContinuousAtZero}
\leanok
\end{lemma}
\begin{proof}\leanok
Identical to the proof of
Lemma~\ref{lem:identity-chart-coeff-contdiff}, applied in the
inversion chart in place of the identity chart. Same chart-
trivialisation Mathlib gap.
\end{proof}

\begin{lemma}[No-\(\infty\) off the pole locus]
\label{lem:meromorphic-no-infty-off-pole}

A meromorphic map \(f:X\to\CPone\) with pole divisor \([P]\) takes
finite values off \(P\): \(\forall x\neq P,\ f(x)\neq\infty\).
\lean{JacobianChallenge.HolomorphicForms.MeromorphicMapToSphere.toMap_ne_infty_off_pole}
\leanok
\end{lemma}
\begin{proof}\leanok
The pole-divisor records the exact \(\infty\)-locus; pole divisor
\([P]\) means only \(P\) is sent to \(\infty\).
\end{proof}

\begin{lemma}[Continuity from no-\(\infty\)]
\label{lem:meromorphic-continuous-on-finite-set}

On any subset \(S\subseteq X\) where \(f\) takes only finite values,
\(f:X\to\CPone\) is continuous on \(S\).
\lean{JacobianChallenge.HolomorphicForms.MeromorphicMapToSphere.continuousOn_of_no_infty_on}
\leanok
\end{lemma}
\begin{proof}\leanok
On \(S\), \(f\) factors through \(\iota:\C\to\CPone\) via a smooth
\(g:S\to\C\); both \(g\) and \(\iota\) are continuous, so the
composition is continuous on \(S\).
\end{proof}

\begin{lemma}[Modulus diverges at a simple pole]
\label{lem:simple-pole-modulus-diverges}

For a meromorphic map \(f:X\to\CPone\) with \(\operatorname{poles}(f)=[P]\), there
exists a smooth \(\C\)-valued lift \(g:X\setminus\{P\}\to\C\)
satisfying \(f|_{X\setminus\{P\}}=\iota\circ g\) and
\(\|g(x)\|\to\infty\) as \(x\to P\).
\lean{JacobianChallenge.HolomorphicForms.MeromorphicMapToSphere.modulus_diverges_at_simple_pole}
\leanok
\end{lemma}
\begin{proof}\leanok
Off \(P\), use Lemma~\ref{lem:meromorphic-no-infty-off-pole} to
construct \(g\). Near \(P\), the local Laurent expansion of \(f\)
takes the form \(z\mapsto a/(z-z_{0})+\text{holomorphic}\) with
\(a\neq 0\) (simple pole), so \(\|g(z)\|\sim|a|/|z-z_{0}|\to\infty\).
\end{proof}

\begin{lemma}[Modulus divergence implies \(\code{OnePoint}\)-tendsto-\(\infty\)]
\label{lem:onepoint-tendsto-infty-of-modulus-diverges}
A \(\C\)-valued function whose modulus diverges along a punctured
neighbourhood of \(P\), lifted to \(\code{OnePoint}\,\C\) via the
canonical inclusion, tends to \(\infty\) at \(P\).
\lean{JacobianChallenge.HolomorphicForms.OnePoint.tendsto_infty_of_modulus_diverges}
\leanok
\end{lemma}
\begin{proof}\leanok
\(\|g\|\to\infty\) along the punctured filter means
\(\code{Tendsto}\,g\) lands in \(\code{Bornology.cobounded}\,\C\) by
\code{Metric.tendsto\_dist\_right\_atTop\_iff} (with base point
\(0\)). On the proper Hausdorff space \(\C\),
\(\code{cobounded}=\code{cocompact}=\code{coclosedCompact}\) by
\code{Metric.cobounded\_eq\_cocompact} and
\code{Filter.coclosedCompact\_eq\_cocompact}. Compose with
\code{OnePoint.tendsto\_coe\_infty} to obtain
\(\iota\circ g\to\infty\) in \(\code{OnePoint}\,\C\).
\end{proof}

% NOTE (audit): lem:degree-one-no-ramification-detail and
% lem:no-ramification-fiber-card-constant entries were removed from
% this section after a cross-section audit found that the existing
% lem:degree-one-no-ramification (Sec02 Blueprint) covers the same
% mathematical content sorry-free at the BranchedCoverData level. The
% MeromorphicDegree.lean route was refactored to use the M-bridge
% (`exists_branchedCoverData_of_pole_degree_one`) and route through
% Blueprint.degree_one_bijective + Blueprint.degree_one_no_ramification
% instead of re-proving the content at the MeromorphicMapToSphere level.

% NOTE (audit): `lem:meromorphic-image-iscompact` and
% `lem:degree-one-open-map` entries were removed from this section
% after the cross-section duplicate audit. They were lemmas in the
% parallel `MeromorphicMapToSphere`-level surjectivity chain
% (M4a + M4b in the project's internal labelling). After M4 was
% refactored to route through the M-bridge to
% `Blueprint.degree_one_bijective`, both lemmas became dead code:
% no reachable node `` either of them. The corresponding Lean
% obligation `image_isCompact` is sorry-free as a one-line wrapper
% around `IsCompact.image`; `isOpenMap_of_pole_degree_one` has been
% removed from `Jacobian/HolomorphicForms/MeromorphicDegree.lean`.

\begin{lemma}[Cotangent transition formula on the overlap]
\label{lem:cotangent-transition-formula}
\uses{lem:onepoint-cx-chart-overlap-pullback}
For \(\omega\in H^{0}(\CPone,\Omega^{1}_{\CPone})\), there is a
punctured neighbourhood of \(0\in\C\) (in the inversion chart) on
which the identity-chart coefficient \(f\) and the inversion-chart
coefficient \(g\) of \(\omega\) are related by
\[
  f(w^{-1})=-w^{2}g(w),\qquad w\neq 0.
\]
\lean{JacobianChallenge.HolomorphicForms.holomorphicOneForm_identityInversionTransition_eventually}
\leanok
\end{lemma}
\begin{proof}\leanok
Pointwise from Lemma~\ref{lem:onepoint-cx-chart-overlap-pullback};
the eventually-quantifier on the punctured neighbourhood follows
via \code{filter\_upwards}.
\end{proof}

\begin{lemma}[Derivative of inverse on the chart overlap]
\label{lem:onepoint-cx-chart-overlap-derivative}
For \(w\neq 0\), \(\frac{d}{dw}(w^{-1})=-w^{-2}\).
\lean{JacobianChallenge.HolomorphicForms.onePointCx_chart_overlap_derivative}
\leanok
\end{lemma}
\begin{proof}\leanok
Direct from Mathlib's \code{hasDerivAt\_inv} on \(w\neq 0\).
\end{proof}

\begin{lemma}[Cotangent pullback on the chart overlap]
\label{lem:onepoint-cx-chart-overlap-pullback}
\uses{lem:onepoint-cx-chart-overlap-derivative}
For \(\omega\in H^{0}(\CPone,\Omega^{1}_{\CPone})\) and \(w\neq 0\),
\(f(w^{-1})=-w^{2}g(w)\), where \(f\) and \(g\) are respectively the
identity-chart and inversion-chart coefficients.
\lean{JacobianChallenge.HolomorphicForms.holomorphicOneForm_chartOverlap_pullback}
\leanok
\end{lemma}
\begin{proof}\leanok
Apply the chain rule for cotangent vectors under
\(z\leftrightarrow w=z^{-1}\):
\(f(z)\,dz=f(w^{-1})\,d(w^{-1})=-f(w^{-1})\,w^{-2}\,dw\). Equating
with \(g(w)\,dw\) yields the identity.
\end{proof}

\begin{lemma}[Finite coefficient decays at infinity]
\label{lem:holomorphic-one-form-coeff-tendsto-zero}
\uses{lem:inversion-chart-coeff-continuous-at-zero, lem:cotangent-transition-formula}
For \(\omega\in H^{0}(\CPone,\Omega^{1}_{\CPone})\), the identity-chart
coefficient tends to \(0\) along the cocompact filter on \(\C\).
\lean{JacobianChallenge.HolomorphicForms.holomorphicOneForm_coeff_tendsto_zero}
\leanok
\end{lemma}
\begin{proof}\leanok
Continuity of the inversion-chart coefficient at \(0\), together with
Lemma~\ref{lem:cotangent-transition-formula}, shows that
\(f(w^{-1})=-w^{2}g(w)\) tends to \(0\) along the punctured neighbourhood of
\(0\). Mathlib's \code{Filter.tendsto\_inv₀\_cobounded'} turns
the inverse map into a map from the cobounded filter to that punctured
neighbourhood, and \code{Metric.cobounded\_eq\_cocompact} identifies the
cobounded and cocompact filters on \(\C\). Hence \(f\to0\) at infinity.
\end{proof}

\begin{lemma}[Liouville vanishing for entire functions decaying at infinity]
\label{lem:entire-tendsto-zero-eq-zero}
An entire function \(f:\C\to\C\) which tends to \(0\) along the
cocompact filter is identically zero.
\lean{JacobianChallenge.HolomorphicForms.entire_tendsto_zero_eq_zero}
\leanok
\end{lemma}
\begin{proof}\leanok
This is Mathlib's Liouville-style lemma
\code{Differentiable.eq\_zero\_of\_tendsto\_zero\_cocompact}, applied to the
given complex differentiable function and its cocompact decay hypothesis.
\end{proof}

\begin{lemma}[Finite-point vanishing on \(\code{OnePoint}\,\C\)]
\label{lem:onepoint-holomorphic-one-form-finite-vanishing}
\uses{lem:entire-tendsto-zero-eq-zero, lem:holomorphic-one-form-coeff-entire, lem:holomorphic-one-form-coeff-tendsto-zero}
For a holomorphic \(1\)-form \(\omega\) on \(\CPone\), the value of
\(\omega\) at every finite point \(z:\C\subset\CPone\) is zero.
\lean{JacobianChallenge.HolomorphicForms.holomorphicOneForm_onePointCx_toFun_finite_eq_zero}
\leanok
\end{lemma}
\begin{proof}\leanok
Read \(\omega\) in the identity chart as \(f(z)\,dz\).
Lemma~\ref{lem:holomorphic-one-form-coeff-entire} supplies the entire coefficient,
while Lemma~\ref{lem:holomorphic-one-form-coeff-tendsto-zero} supplies
cocompact decay of \(f\). Applying
Lemma~\ref{lem:entire-tendsto-zero-eq-zero} gives \(f=0\), hence \(\omega\)
vanishes at all finite points.
\end{proof}

\begin{lemma}[Pointwise vanishing on \(\code{OnePoint}\,\C\)]
\label{lem:onepoint-holomorphic-one-form-pointwise-vanishing}
\uses{lem:onepoint-holomorphic-one-form-finite-vanishing, lem:inversion-chart-coeff-continuous-at-zero}
Every holomorphic \(1\)-form on \(\CPone\) evaluates to zero at every point,
including the point at infinity.
\lean{JacobianChallenge.HolomorphicForms.holomorphicOneForm_onePointCx_toFun_eq_zero}
\leanok
\end{lemma}
\begin{proof}\leanok
Finite points are covered by
Lemma~\ref{lem:onepoint-holomorphic-one-form-finite-vanishing}. At infinity,
the inversion-chart coefficient is continuous at \(0\) by
Lemma~\ref{lem:inversion-chart-coeff-continuous-at-zero}; since it vanishes on
the punctured neighbourhood by finite-point vanishing, continuity forces the
value at \(0\), i.e. the value at \(\infty\), to vanish as well.
\end{proof}

\begin{lemma}[Holomorphic \(1\)-forms on \(\code{OnePoint}\,\C\) are subsingletons]
\label{lem:onepoint-holomorphic-one-form-subsingleton}
\uses{lem:onepoint-holomorphic-one-form-pointwise-vanishing}
The space \(H^{0}(\CPone,\Omega^{1}_{\CPone})\) is a subsingleton.
\lean{JacobianChallenge.HolomorphicForms.holomorphicOneForm_onePointCx_subsingleton}
\leanok
\end{lemma}
\begin{proof}\leanok
Two holomorphic \(1\)-forms on \(\CPone\) have the same pointwise value, since
both vanish everywhere by
Lemma~\ref{lem:onepoint-holomorphic-one-form-pointwise-vanishing}. Hence the
forms are equal.
\end{proof}

\begin{lemma}[Analytic genus of \(\code{OnePoint}\,\C\)]
\label{lem:analytic-genus-onepoint-zero}
\uses{lem:onepoint-holomorphic-one-form-subsingleton}
The analytic genus of \(\CPone=\code{OnePoint}\,\C\) is zero.
\lean{JacobianChallenge.HolomorphicForms.analyticGenus_onePointCx_eq_zero}
\leanok
\end{lemma}
\begin{proof}\leanok
By Lemma~\ref{lem:onepoint-holomorphic-one-form-subsingleton}, the module of
holomorphic \(1\)-forms on \(\CPone\) is a subsingleton, hence zero-dimensional.
The definition of analytic genus is this complex dimension.
\end{proof}

\begin{lemma}[Analytic genus transported from a sphere homeomorphism]
\label{lem:analytic-genus-homeomorphic-sphere-eq-onepoint}
\uses{lem:holomorphic-one-form-equiv-of-homeo-sphere}
If a compact connected Riemann surface \(X\) is homeomorphic to the standard
\(2\)-sphere, then its analytic genus equals that of
\(\code{OnePoint}\,\C\).
\lean{JacobianChallenge.HolomorphicForms.analyticGenus_eq_of_homeomorphic_sphere_of_onePointCx}
Lemma~\ref{lem:holomorphic-one-form-equiv-of-homeo-sphere} gives a
\(\C\)-linear equivalence between the holomorphic \(1\)-form spaces on \(X\)
and on \(\code{OnePoint}\,\C\). Analytic genus is the complex dimension of this
space, so the finite ranks are equal.
\end{lemma}

\begin{theorem}[Analytic genus vanishes for surfaces homeomorphic to \(S^2\)]
\label{thm:analytic-genus-zero-of-homeomorphic-sphere}
\uses{lem:analytic-genus-homeomorphic-sphere-eq-onepoint, lem:analytic-genus-onepoint-zero}
A compact connected Riemann surface homeomorphic to the standard \(2\)-sphere
has analytic genus zero.
\lean{JacobianChallenge.HolomorphicForms.analyticGenus_eq_zero_of_homeomorphic_sphere}
By Lemma~\ref{lem:analytic-genus-homeomorphic-sphere-eq-onepoint}, the analytic
genus of \(X\) is the analytic genus of \(\code{OnePoint}\,\C\). That value is
zero by Lemma~\ref{lem:analytic-genus-onepoint-zero}.
\end{theorem}

\subsection{Sup-norm completeness decomposition (CompactRiemannSurface)}
\label{sec:sup-norm-completeness-decomp}

\begin{lemma}[Pointwise limit exists]
\label{lem:holomorphic-one-form-pointwise-limit-exists}
A sup-norm Cauchy sequence of holomorphic 1-forms on a compact
Riemann surface admits a pointwise limit at every point of \(X\).
\lean{JacobianChallenge.HolomorphicForms.holomorphicOneForm_supNorm_cauchySeq_pointwise_limit_exists}
\leanok
\end{lemma}
\begin{proof}\leanok
Sup-norm Cauchy implies pointwise Cauchy in each fiber
\(\mathrm{Cot}_x(X)\), which is a Banach space, so the pointwise limit
exists.
\end{proof}

\begin{lemma}[Continuity of the pointwise limit]
\label{lem:holomorphic-one-form-limit-continuous}
\uses{lem:holomorphic-one-form-pointwise-limit-exists}
The pointwise limit of a sup-norm Cauchy sequence of holomorphic
1-forms is continuous.
\lean{JacobianChallenge.HolomorphicForms.holomorphicOneForm_supNorm_cauchySeq_limit_continuous}
\leanok
\end{lemma}
\begin{proof}\leanok
Sup-norm Cauchy convergence is uniform; the uniform limit of
continuous functions is continuous (\code{TendstoUniformly.continuous}).
\end{proof}

\begin{lemma}[Holomorphicity of the pointwise limit (Weierstrass)]
\label{lem:holomorphic-one-form-limit-holomorphic}
\uses{lem:holomorphic-one-form-pointwise-limit-exists}
The pointwise limit of a sup-norm Cauchy sequence of holomorphic
1-forms is holomorphic.
\lean{JacobianChallenge.HolomorphicForms.holomorphicOneForm_supNorm_cauchySeq_limit_holomorphic}
\leanok
\end{lemma}
\begin{proof}\leanok
Classical Weierstrass convergence theorem: a uniform limit of
holomorphic functions is holomorphic, via Morera's theorem (closed
contour integrals vanish under uniform convergence) or
power-series-coefficient convergence (Cauchy estimates pass to the
limit). \textbf{Mathlib gap}: v4.28.0 lacks a Weierstrass-convergence
theorem for holomorphic functions or sections.
\end{proof}

\subsection{Uniformization-lite further decomposition}
\label{sec:uniformization-lite-iter3}

\begin{lemma}[Homeomorphism to \(\code{OnePoint}\,\C\) from a sphere homeomorphism]
\label{lem:uniformization-genus-zero-homeomorphism}
A compact connected complex 1-manifold homeomorphic to \(S^2\) admits a
\emph{homeomorphism} to \(\code{OnePoint}\,\C\).
\lean{JacobianChallenge.HolomorphicForms.exists_biholomorphism_to_OnePointCx_of_homeoSphere}
\leanok
\end{lemma}
\begin{proof}
\leanok
Purely topological: compose the given \(X\simeq_{\mathrm{top}}S^2\) with the
fixed homeomorphism \(S^2\simeq_{\mathrm{top}}\code{OnePoint}\,\C\). This yields
only a homeomorphism; the complex-smooth upgrade is
Theorem~\ref{thm:uniformization-genus-zero-biholomorphism}.
\end{proof}

\begin{lemma}[Gluing coordinates land in their target charts]
\label{lem:global-gluing-coord-mem-target}
For a genus-zero global patch family, each local coordinate value lies in the
target of the assigned public \(\code{OnePoint}\,\C\) chart.
\lean{JacobianChallenge.HolomorphicForms.genusZeroGlobalGluing_coord_mem_target_on_patch}
\leanok
\end{lemma}
\begin{proof}\leanok
The assigned target chart is either the identity chart or the inversion chart
of \(\code{OnePoint}\,\C\). In both cases the coordinate target-membership
condition follows by unfolding the standard chart target.
\end{proof}

\begin{definition}[Glued forward candidate map]
\label{def:global-gluing-to-map}
\uses{lem:global-gluing-coord-mem-target}
Given a genus-zero global patch family, define the candidate map
\(X\to\code{OnePoint}\,\C\) by choosing a patch containing \(x\) and applying
that patch's target-chart inverse to the local coordinate of \(x\).
\lean{JacobianChallenge.HolomorphicForms.genusZeroGlobalGluing_toMap}
\leanok
\end{definition}

\begin{lemma}[Glued map agrees with the selector uniformization]
\label{lem:global-gluing-to-map-eq-uniformization}
\uses{def:global-gluing-to-map}
For a normalized Montel patch selector, the chosen-patch formula for the glued
map agrees pointwise with the homeomorphism packaged by the selector.
\lean{JacobianChallenge.HolomorphicForms.genusZeroGlobalGluing_toMap_eq_uniformization}
\leanok
\end{lemma}
\begin{proof}\leanok
Choose the patch supplied by the finite cover at \(x\). The selector field
\code{coord\_represents\_uniformization} states exactly that this patch's
target-chart inverse of the local coordinate is the selector uniformization.
\end{proof}

\subsection{Montel coordinate representation}
\label{sec:montel-coord-represents}

This subtree isolates the \code{coord\_represents\_uniformization} field of
\code{GenusZeroNormalizedMontelPatchSelector}.  It is the local statement that
each selected Montel coordinate represents the global uniformization in the
assigned public chart of \(\code{OnePoint}\,\C\).

\begin{lemma}[Montel selector coordinate represents the uniformization (open)]
\label{lem:montel-selector-coord-represents-uniformization}
\uses{lem:montel-selector-source-patch-cover, lem:global-gluing-coord-mem-target, lem:montel-selector-chart-inverse-local-value, lem:montel-selector-normal-limit-agrees-with-uniformization, lem:montel-selector-overlap-compatible-normalized-values}
For every selected source patch \(U_i\), source point \(x\in U_i\), and assigned
target chart \(\psi_i\) on \(\code{OnePoint}\,\C\), the selected local coordinate
represents the global uniformization:
\[
  \psi_i^{-1}(z_i(x)) = F(x).
\]
This is the field
\code{GenusZeroNormalizedMontelPatchSelector.coord\_represents\_uniformization}.
The Lean projection exists as a structure field, but the construction of a
selector carrying this field is the open Montel assembly obligation, so this
blueprint node intentionally carries no \code{\textbackslash lean\{\}} citation.
\end{lemma}

\begin{lemma}[Selected source patches cover the surface]
\label{lem:montel-selector-source-patch-cover}
The finite global patch-family structure supplies a selected source patch
containing each point of \(X\).  Its \code{patch\_cover} field is the cover
witness used when turning local coordinate statements into a global formula.
\lean{JacobianChallenge.HolomorphicForms.GenusZeroGlobalPatchFamily}
\leanok
\end{lemma}
\begin{proof}\leanok\end{proof}

\begin{lemma}[Selected chart inverse gives the local candidate value (open)]
\label{lem:montel-selector-chart-inverse-local-value}
\uses{lem:global-gluing-coord-mem-target, lem:mathlib-open-partial-homeomorph-map-target, lem:mathlib-open-partial-homeomorph-right-inv}
On a selected source patch, the local coordinate lies in the assigned target
chart.  Applying the target-chart inverse therefore gives the local candidate
point of \(\code{OnePoint}\,\C\) that should become the global uniformization
value.

This is a planning node for the project-specific reconstruction package.  Its
formal floor is already visible: target membership comes from the global patch
family, and Mathlib's open-partial-homeomorphism API says that the chart inverse
lands in the chart source and reconstructs the original coordinate.
\end{lemma}

\begin{lemma}[Mathlib: target points map back to the source]
\label{lem:mathlib-open-partial-homeomorph-map-target}
For an open partial homeomorphism \(e\), if \(z\in e.\mathrm{target}\), then
\(e^{-1}(z)\in e.\mathrm{source}\).
\lean{OpenPartialHomeomorph.map_target}
\leanok
\end{lemma}
\begin{proof}\leanok\end{proof}

\begin{lemma}[Mathlib: chart inverse reconstructs target coordinates]
\label{lem:mathlib-open-partial-homeomorph-right-inv}
For an open partial homeomorphism \(e\), if \(z\in e.\mathrm{target}\), then
\[
  e(e^{-1}(z)) = z.
\]
\lean{OpenPartialHomeomorph.right_inv}
\leanok
\end{lemma}
\begin{proof}\leanok\end{proof}

\begin{lemma}[Normalized Montel limit agrees with the uniformization (open)]
\label{lem:montel-selector-normal-limit-agrees-with-uniformization}
\uses{lem:montel-selector-source-patch-cover, lem:montel-selector-chart-inverse-local-value}
The normal-family/Montel extraction must choose local coordinates whose
target-chart inverse agrees with the global homeomorphism \(F:X\simeq
\code{OnePoint}\,\C\) on the selected source patch.

This is the analytic heart of the coordinate-representation field: it is not a
formal consequence of the existing gluing API, and it should remain an open
provider until the Montel patch-selector construction supplies it.
\end{lemma}

\begin{lemma}[Normalized local values are overlap-compatible (open)]
\label{lem:montel-selector-overlap-compatible-normalized-values}
\uses{lem:global-gluing-coord-mem-target}
On overlaps of selected source patches, the two target-chart inverse formulas
for the normalized Montel coordinates determine the same point of
\(\code{OnePoint}\,\C\).  This compatibility is what lets the local
representations define one global uniformization rather than unrelated chart
values.

This is a planning node with no current Lean declaration.  Downstream green
gluing lemmas consume this compatibility once the selector field is available,
but they should not be used as fake providers for constructing the field.
\end{lemma}


\begin{lemma}[Gluing coordinates agree on overlaps]
\label{lem:global-gluing-overlap-compatible}
\uses{lem:global-gluing-to-map-eq-uniformization}
On any overlap of two selected source patches, the two local coordinate
formulas determine the same point of \(\code{OnePoint}\,\C\).
\lean{JacobianChallenge.HolomorphicForms.genusZeroGlobalGluing_overlap_compatible}
\leanok
\end{lemma}
\begin{proof}\leanok
Both patch expressions are equal to the selector's uniformization at the same
source point, so they are equal to each other.
\end{proof}

\begin{lemma}[Glued map lands in the selected target chart]
\label{lem:global-gluing-to-map-target-mem}
\uses{lem:global-gluing-coord-mem-target, lem:global-gluing-to-map-eq-uniformization}
On each selected source patch, the glued map lands in the source of that
patch's assigned target chart.
\lean{JacobianChallenge.HolomorphicForms.genusZeroGlobalGluing_toMap_target_mem_on_patch}
\leanok
\end{lemma}
\begin{proof}\leanok
The local coordinate lies in the target chart by
Lemma~\ref{lem:global-gluing-coord-mem-target}. Applying the chart inverse
puts the value in the chart source, and
Lemma~\ref{lem:global-gluing-to-map-eq-uniformization} identifies this local
formula with the chosen-patch glued map.
\end{proof}

\begin{lemma}[Forward glued candidate exists]
\label{lem:global-gluing-to-map-exists}
\uses{def:global-gluing-to-map, lem:global-gluing-to-map-target-mem}
A normalized Montel patch selector supplies a global candidate
\(X\to\code{OnePoint}\,\C\) whose values lie in the assigned target chart on
each source patch.
\lean{JacobianChallenge.HolomorphicForms.genusZeroGlobalGluing_toMap_exists}
\leanok
\end{lemma}
\begin{proof}\leanok
Take the map of Definition~\ref{def:global-gluing-to-map}; the target-chart
membership assertion is Lemma~\ref{lem:global-gluing-to-map-target-mem}.
\end{proof}

\begin{lemma}[Chart expression for the glued map]
\label{lem:global-gluing-chart-expression}
\uses{lem:global-gluing-coord-mem-target, lem:global-gluing-to-map-eq-uniformization}
On each selected source patch, reading the glued map in the assigned target
chart recovers the patch's local coordinate.
\lean{JacobianChallenge.HolomorphicForms.genusZeroGlobalGluing_chart_expression_on_patch}
\leanok
\end{lemma}
\begin{proof}\leanok
After replacing the glued map by the selector uniformization using
Lemma~\ref{lem:global-gluing-to-map-eq-uniformization}, the selector's local
coordinate representation reduces the claim to the right-inverse property of
the target chart.
\end{proof}

\begin{lemma}[Glued inverse candidate exists]
\label{lem:global-gluing-inv-map-exists}
\uses{lem:global-gluing-chart-expression}
The selector supplies an inverse candidate \(\code{OnePoint}\,\C\to X\) whose
local formula on each target chart is the selected inverse branch.
\lean{JacobianChallenge.HolomorphicForms.genusZeroGlobalGluing_invMap_exists}
\leanok
\end{lemma}
\begin{proof}\leanok
Use the inverse of the selector homeomorphism. The selector field
\code{invCoord\_represents\_uniformization} identifies this inverse with each
local inverse branch on the corresponding target chart.
\end{proof}

\begin{lemma}[Local left inverse for the glued maps]
\label{lem:global-gluing-local-left-inverse}
\uses{lem:global-gluing-to-map-eq-uniformization}
The selector inverse followed by the glued forward map is locally the identity
on every selected source patch.
\lean{JacobianChallenge.HolomorphicForms.genusZeroGlobalGluing_local_left_inverse_on_patch}
\leanok
\end{lemma}
\begin{proof}\leanok
Replace the glued forward map by the selector uniformization, then apply the
left-inverse property of the selector homeomorphism.
\end{proof}

\begin{lemma}[Local right inverse for target-chart branches]
\label{lem:global-gluing-local-right-inverse}
\uses{lem:global-gluing-to-map-eq-uniformization, lem:global-gluing-inv-map-exists}
The glued forward map sends each selected inverse branch back to the
corresponding point of \(\code{OnePoint}\,\C\).
\lean{JacobianChallenge.HolomorphicForms.genusZeroGlobalGluing_local_right_inverse_on_target_chart}
\leanok
\end{lemma}
\begin{proof}\leanok
Rewrite the glued map as the selector uniformization and rewrite the inverse
branch by \code{invCoord\_represents\_uniformization}; the result is the
right-inverse property of the selector homeomorphism.
\end{proof}

\begin{lemma}[Smoothness of the glued forward map]
\label{lem:global-gluing-contmdiff-to-map}
\uses{lem:global-gluing-to-map-eq-uniformization}
The glued forward candidate is \(\code{ContMDiff}\).
\lean{JacobianChallenge.HolomorphicForms.genusZeroGlobalGluing_contMDiff_toMap}
\leanok
\end{lemma}
\begin{proof}\leanok
The glued map is pointwise the selector uniformization by
Lemma~\ref{lem:global-gluing-to-map-eq-uniformization}; smoothness is the
\code{uniformization\_contMDiff} field of the selector.
\end{proof}

\begin{lemma}[Smoothness of the glued inverse map]
\label{lem:global-gluing-contmdiff-inv-map}
\uses{lem:global-gluing-inv-map-exists}
The inverse map packaged by the normalized Montel selector is
\(\code{ContMDiff}\).
\lean{JacobianChallenge.HolomorphicForms.genusZeroGlobalGluing_contMDiff_invMap}
\leanok
\end{lemma}
\begin{proof}\leanok
This is exactly the \code{inverse\_uniformization\_contMDiff} field of the
normalized Montel patch selector.
\end{proof}

\begin{lemma}[Montel chart-ball equicontinuity]
\label{lem:montel-chartball-equicontinuity}
Uniform Cauchy bounds on a family of packaged chart-ball functions give the
uniform epsilon-delta equicontinuity estimate on the smaller chart ball.
\lean{JacobianChallenge.HolomorphicForms.ChartBallPowerSeries.family_uniform_equicontinuousOn_of_cauchy_bound}
\leanok
\end{lemma}

\begin{lemma}[Montel chart-ball Arzel\`a--Ascoli compact closure]
\label{lem:montel-chartball-arzela-ascoli-compact}
\uses{lem:montel-chartball-equicontinuity}
After restriction to a compact closed ball, the equicontinuous chart-ball
family has compact closure in the bounded-continuous function space by
Mathlib's Arzel\`a--Ascoli theorem.
\lean{JacobianChallenge.HolomorphicForms.ChartBallPowerSeries.isCompact_closure_boundedContinuousOnClosedBall_range}
\leanok
\end{lemma}

\begin{lemma}[Montel chart-ball subsequence extraction]
\label{lem:montel-chartball-subsequence-extraction}
\uses{lem:montel-chartball-arzela-ascoli-compact}
The compact-closure statement extracts a subsequence converging in the
bounded-continuous closed-ball function space.
\lean{JacobianChallenge.HolomorphicForms.ChartBallPowerSeries.exists_tendsto_subseq_boundedContinuousOnClosedBall}
\leanok
\end{lemma}

\begin{lemma}[Montel chart-ball uniform closed-ball limit]
\label{lem:montel-chartball-uniform-closed-ball-limit}
\uses{lem:montel-chartball-subsequence-extraction}
The closed-ball bounded-continuous subsequence convergence projects back to
uniform convergence of the original chart functions on that compact closed
ball.
\lean{JacobianChallenge.HolomorphicForms.ChartBallPowerSeries.exists_subseq_tendstoUniformlyOn_closedBall}
\leanok
\end{lemma}

\begin{lemma}[Montel finite-family diagonal extraction]
\label{lem:montel-chartball-finite-family-diagonal-extraction}
\uses{lem:montel-chartball-uniform-closed-ball-limit}
Given a finite family of equicontinuous chart-ball power series sequences, 
there exists a single strictly monotone subsequence of indices along which
all sequences converge uniformly on their respective closed balls simultaneously.
\lean{JacobianChallenge.HolomorphicForms.exists_subseq_tendstoUniformlyOn_closedBall_finite}
\leanok
\end{lemma}

\begin{lemma}[Montel finite-family locally uniform diagonal extraction]
\label{lem:montel-finite-family-locally-uniform-diagonal-extraction}
\uses{lem:montel-chartball-finite-family-diagonal-extraction}
Given a finite family of equicontinuous chart-ball power series sequences on nested subballs covering an open ball, there exists a single strictly monotone subsequence of indices along which all sequences converge locally uniformly on their respective open balls simultaneously.
\lean{JacobianChallenge.HolomorphicForms.exists_diagonal_subseq_tendstoLocallyUniformlyOn_finite}
\leanok
\end{lemma}

\begin{lemma}[Montel chart-ball Weierstrass limit]
\label{lem:montel-chartball-weierstrass-limit}
\uses{lem:montel-chartball-uniform-closed-ball-limit}
A locally uniform chart-ball limit is differentiable on the common open chart
ball, and the derivatives converge locally uniformly to the derivative of the
limit.
\lean{JacobianChallenge.HolomorphicForms.ChartBallPowerSeries.differentiableOn_limit_of_tendstoLocallyUniformlyOn_chartBall, JacobianChallenge.HolomorphicForms.ChartBallPowerSeries.tendstoLocallyUniformlyOn_deriv_of_chartBall_limit}
\leanok
\end{lemma}

\begin{lemma}[Montel chart-ball normalized local chart]
\label{lem:montel-chartball-normalized-local-chart}
\uses{lem:montel-chartball-weierstrass-limit}
Locally uniform convergence preserves the point and nonzero-derivative
normalization, and the resulting normalized chart-ball limit promotes to the
local chart-homeomorphism data used by the genus-zero gluing construction.
\lean{JacobianChallenge.HolomorphicForms.ChartBallPowerSeries.normalizedChartBallLimit_of_tendstoLocallyUniformlyOn, JacobianChallenge.HolomorphicForms.ChartBallPowerSeries.exists_localNormalizedChartHomeomorphData_of_tendstoLocallyUniformlyOn}
\leanok
\end{lemma}

\begin{lemma}[Montel local chart patch from a normalized chart-ball limit]
\label{lem:montel-local-chart-patch-from-limit}
\uses{lem:montel-chartball-normalized-local-chart}
A packaged normalized chart-ball limit, together with its local
inverse-function-theorem chart and a choice of public \(\code{OnePoint}\,\C\)
target chart, gives the local Montel chart patch data.
\lean{JacobianChallenge.HolomorphicForms.GenusZeroLocalMontelChartPatch.ofChartBallLimit}
\leanok
\end{lemma}

\begin{lemma}[Montel local patch realization (open)]
\label{lem:montel-local-patch-realization}
\uses{lem:montel-local-chart-patch-from-limit}
Each selected global gluing patch must be realized by its local normalized
Montel chart: the source chart lands in the chart-ball domain, the patch
coordinate is the chart-ball limit in that source coordinate, and the inverse
branch agrees with the local inverse. This is the real
\code{GenusZeroLocalMontelPatchRealization} payload, but no green existence
provider currently selects these realizations for the finite family.
\lean{JacobianChallenge.HolomorphicForms.GenusZeroLocalMontelPatchRealization}
\end{lemma}

\begin{lemma}[Chart-ball map agrees with the local chart on its source]
\label{lem:montel-realization-toFun-eq-localOpen}
The packaged chart-ball map \(\code{chartBall.toFun}\) agrees, on the source of
the local inverse-function-theorem chart, with the topological partial
homeomorphism \(\code{localOpen}\). This is the analytic/topological bridge used
to express the realized patch inverse branch; it is the
\code{toFun\_eq\_on\_source} field of the local chart data read off as an
equality of maps.
\lean{JacobianChallenge.HolomorphicForms.ChartBallPowerSeries.toFun_eq_localOpen_on_source}
\leanok
\end{lemma}
\begin{proof}\leanok\end{proof}

\begin{lemma}[Realized patch local right-inverse]
\label{lem:montel-realization-coord-invcoord}
\uses{lem:montel-realization-toFun-eq-localOpen, lem:mathlib-open-partial-homeomorph-right-inv}
For the patch built by the green local-patch realization provider, on the local
chart target the patch coordinate undoes its inverse branch:
\(\code{coord}(\code{invCoord}\,z)=z\). This is the local right-inverse fact the
global-gluing substrate consumes; it follows from the section identity, the
chart-ball/local-chart agreement, and \code{OpenPartialHomeomorph.right\_inv}.
\lean{JacobianChallenge.HolomorphicForms.montelRealizationPatch_coord_invCoord}
\leanok
\end{lemma}
\begin{proof}\leanok\end{proof}

\begin{lemma}[Realized patch local left-inverse]
\label{lem:montel-realization-invcoord-coord}
\uses{lem:montel-realization-toFun-eq-localOpen, lem:mathlib-open-partial-homeomorph-left-inv}
For the patch built by the green local-patch realization provider, on the patch
source the inverse branch undoes the coordinate:
\(\code{invCoord}(\code{coord}\,x)=x\). This is the local left-inverse fact the
global-gluing substrate consumes; it follows from the chart-ball/local-chart
agreement, \code{OpenPartialHomeomorph.left\_inv} on the domain ball, and one
honest section round-trip hypothesis a caller supplies.
\lean{JacobianChallenge.HolomorphicForms.montelRealizationPatch_invCoord_coord}
\leanok
\end{lemma}
\begin{proof}\leanok\end{proof}

\begin{lemma}[Packaged realized Montel patch bundle]
\label{lem:montel-realized-patch-bundle}
\uses{lem:montel-local-patch-realization, lem:montel-realization-coord-invcoord, lem:montel-realization-invcoord-coord}
A single bundle for the downstream gluing consumers: the constructed global
gluing patch, its local Montel realization witness, and both local round-trip
identities (\(\code{coord}\circ\code{invCoord}=\mathrm{id}\) on the local chart
target, \(\code{invCoord}\circ\code{coord}=\mathrm{id}\) on the patch source). The
constructor assembles it from the chart-ball substrate and source-chart data as a
pure composition of the three accepted green leaves, so a consumer obtains a
complete, ready-to-glue realized patch in one value rather than re-threading the
source-chart hypotheses across four separate calls.
\lean{JacobianChallenge.HolomorphicForms.MontelRealizedPatch, JacobianChallenge.HolomorphicForms.montelRealizedPatch_of_sourceChart}
\leanok
\end{lemma}
\begin{proof}\leanok\end{proof}

\begin{lemma}[Realized-patch bundle from an honest two-sided section]
\label{lem:montel-realized-patch-section-form}
\uses{lem:montel-realized-patch-bundle}
A convenience form of the realized-patch bundle constructor for the case a gluing
consumer actually instantiates: the section is a genuine two-sided inverse of the
source chart — a global right-inverse
(\(\code{sourceChart}\circ\code{sourceSection}=\mathrm{id}\)) together with a
left-inverse on the patch source. This collapses the general constructor's three
restricted section hypotheses into these two honest facts; the restricted ones are
recovered by instantiating the global right-inverse at
\(w=\code{localOpen.symm}\,z\). Pure delegation to the general constructor.
\lean{JacobianChallenge.HolomorphicForms.montelRealizedPatch_of_sourceChartSection}
\leanok
\end{lemma}
\begin{proof}\leanok\end{proof}

\begin{lemma}[Realized-patch bundle from a domain-restricted section]
\label{lem:montel-realized-patch-local-section-form}
\uses{lem:montel-realized-patch-bundle, lem:montel-realized-patch-section-form}
The most-satisfiable convenience form: the section right-inverse is required only
at the values the realization and round-trip proofs actually touch —
\(\code{localOpen.symm}\,z\) for \(z\in\code{targetChart.target}\cup
\code{localOpen.target}\) — rather than globally on all of \(\C\). The global
two-sided-section form forces \(\code{sourceChart}\) to be surjective onto \(\C\),
which a chart that is a homeomorphism onto a proper open subset cannot satisfy;
the restricted form is dischargeable by such a chart, and the global form factors
through it. Pure delegation to the general constructor via a case split on the
target-membership disjunction.
\lean{JacobianChallenge.HolomorphicForms.montelRealizedPatch_of_sourceChartLocalSection}
\leanok
\end{lemma}
\begin{proof}\leanok\end{proof}

\begin{lemma}[Source-chart package per chart-ball datum]
\label{lem:montel-source-chart-package}
\uses{lem:montel-chartball-normalized-local-chart, lem:montel-realized-patch-local-section-form, lem:mathlib-open-partial-homeomorph-right-inv, lem:mathlib-open-partial-homeomorph-left-inv}
The caller package the finite-cover construction needs, stated abstractly per
chart-ball datum (the selected patches do not exist yet). Given one local
normalized Montel chart patch together with a source chart on \(X\) modeled as an
\code{OpenPartialHomeomorph}\,\(X\,\C\) that realizes its chart-ball coordinate —
smooth on its source, landing inside the chart-ball domain ball, and with the
inverse-function-theorem inverse values of the chart-ball lying in the source
chart's image — it produces a source set, source chart and local section
satisfying \emph{every} hypothesis of the domain-restricted realized-patch
constructor. The witnesses are \(\code{source}=\varphi.\code{source}\),
\(\code{sourceChart}=\varphi\), \(\code{sourceSection}=\varphi.\code{symm}\); the
radius-ball membership follows from the domain ball via
\code{domainRadius\_lt\_chart} and \code{Metric.ball\_subset\_ball}, the section
inverses from \code{OpenPartialHomeomorph.right\_inv}/\code{left\_inv}. The rider
\code{exists\_montelRealizedPatch\_of\_chartBallData} composes this with the
constructor to yield a \code{MontelRealizedPatch} per selected patch in one call.
This is pure packaging plus local chart topology — no new analytic content.
\lean{JacobianChallenge.HolomorphicForms.exists_sourceChartPackage_of_chartBallData}
\leanok
\end{lemma}
\begin{proof}\leanok\end{proof}

\begin{lemma}[Raw-source-chart convenience form of the patch package]
\label{lem:montel-raw-source-chart-package}
\uses{lem:montel-source-chart-package, lem:mathlib-open-partial-homeomorph-right-inv, lem:mathlib-open-partial-homeomorph-left-inv}
The domain-ball-restriction convenience form of the source-chart package, for the
engine's per-patch instantiation point. Starting from a \emph{raw} atlas source
chart \(\varphi:\code{OpenPartialHomeomorph}\,X\,\C\), restrict it via
\code{OpenPartialHomeomorph.restrOpen} to the open set
\(\varphi.\code{source}\cap\varphi^{-1}(\code{ball}\,\code{center}\,\code{domainRadius})\)
(open by \code{isOpen\_inter\_preimage} on the metric ball). On the restricted
chart \(\psi\), the awkward all-of-source image hypothesis \code{himage}
\emph{self-discharges}: every \(x\in\psi.\code{source}\) lies in the preimage, so
\(\psi\,x=\varphi\,x\) is in the domain ball by membership unfolding. The
coercion and inverse transfer \code{rfl}-grade (\code{coe\_restrOpen},
\code{coe\_restrOpen\_symm}), so smoothness transfers by \code{ContMDiffOn.mono}
and the section laws verbatim.

\emph{Honesty note.} The restricted target \(\psi.\code{target}
=\varphi.\code{target}\cap\varphi.\code{symm}^{-1}(\dots)\) is genuinely smaller,
so one might fear an extra \emph{restricted}-image obligation on the
inverse-function-theorem values \(\code{localOpen.symm}\,z\). But the package
conclusion's right-inverse clause is stated about the section
\(\psi.\code{symm}=\varphi.\code{symm}\) applied \emph{directly} to
\(\code{localOpen.symm}\,z\), and \code{OpenPartialHomeomorph.right\_inv}
discharges it from \(\code{localOpen.symm}\,z\in\varphi.\code{target}\) alone —
it never touches \(\psi.\code{target}\). The convenience form is therefore
\emph{strictly weaker} in hypotheses than the chart-ball form: it drops
\code{himage} and keeps only the original \code{hsymm\_image}; no restricted-ball
analogue is owed. The rider
\code{exists\_montelRealizedPatch\_of\_rawChart} composes this with the
local-section constructor to yield a \code{MontelRealizedPatch} from a raw chart
in one call. Pure packaging plus local chart topology — no new analytic content.
\lean{JacobianChallenge.HolomorphicForms.exists_sourceChartPackage_of_rawChart, JacobianChallenge.HolomorphicForms.exists_montelRealizedPatch_of_rawChart}
\leanok
\end{lemma}
\begin{proof}\leanok\end{proof}

\begin{lemma}[Glued-map local agreement on a realized patch family]
\label{lem:montel-glued-map-local-agreement}
\uses{lem:montel-realized-patch-bundle, lem:montel-coherent-local-selector, lem:mathlib-open-partial-homeomorph-right-inv}
The first consumer of the realized-patch toolkit: a selector-free local-formula
step underneath the glued forward map's smoothness. The canonical forward
candidate \(\code{genusZeroGlobalGluing\_toMap}\) is defined for a bare patch
family (\(\code{Classical.choose}\) over the cover, then
\(\code{targetChart.symm}\circ\code{coord}\)); on each source patch \(i\),
applying that patch's target chart to it recovers the patch's own coordinate,
\(\code{targetChart}_i(\code{toMap}\,x)=\code{coord}_i(x)\). Proved from overlap
compatibility \texttt{hcompat} — exactly the conclusion shape of the coherent
local selector's overlap-agreement kernel — and per-patch target membership
\texttt{hmem}, via \(\code{OpenPartialHomeomorph.right\_inv}\), \emph{without}
\(\code{genusZeroGlobalGluing\_toMap\_eq\_uniformization}\) or any normalized
Montel patch selector accessor. The chart-ball corollary then reads the glued
map as the Montel limit in the realized source coordinate, consuming a
\code{MontelRealizedPatch} via its \code{coord\_eq\_chartBall} field.
\lean{JacobianChallenge.HolomorphicForms.genusZeroGlobalGluing_targetChart_toMap_eq_coord, JacobianChallenge.HolomorphicForms.genusZeroGlobalGluing_targetChart_toMap_eq_chartBall}
\leanok
\end{lemma}
\begin{proof}\leanok\end{proof}

\begin{lemma}[Mathlib: atlas-member inverse is smooth]
\label{lem:mathlib-contmdiffon-symm-of-mem-maximalatlas}
For an open partial homeomorphism \(e\) in the maximal atlas of a \(C^n\)
manifold, the inverse \(e^{-1}\) is \(\code{ContMDiffOn}\) on \(e.\mathrm{target}\).
Green (Mathlib).
\lean{contMDiffOn_symm_of_mem_maximalAtlas}
\leanok
\end{lemma}
\begin{proof}\leanok\end{proof}

\begin{lemma}[Per-patch forward smoothness of the candidate uniformization]
\label{lem:montel-forward-per-patch-smoothness}
\uses{lem:montel-glued-map-local-agreement, lem:montel-realized-patch-bundle, lem:mathlib-contmdiffon-symm-of-mem-maximalatlas}
The per-patch smoothness reading — the last local ingredient underneath the
forward-smoothness leaf
(\(\code{genusZeroMontel\_raw\_global\_patch\_family\_contMDiff\_toMap}\)) before
the local-to-global gluing argument. If a Montel patch family represents a
candidate \(u:X\to\code{OnePoint}\,\C\) in local public target charts
(\texttt{hcoord}) and each patch coordinate lands in its target chart's target
(\texttt{hmem}), then \(u\) is \(\code{ContMDiffOn}\) on each patch source. On the
patch, \(u=\code{targetChart.symm}\circ\code{coord}\); \(\code{coord}\) is
\(\code{ContMDiffOn}\) (a patch field) and the standard \(\code{OnePoint}\,\C\)
target chart inverses are \(\code{ContMDiffOn}\) as atlas-member inverses
(\code{targetChart\_symm\_contMDiffOn}, from
Lemma~\ref{lem:mathlib-contmdiffon-symm-of-mem-maximalatlas}). The chart-ball
corollary consumes a \code{MontelRealizedPatch} to identify the source coordinate
with the Montel chart-ball limit. Selector-free.
\lean{JacobianChallenge.HolomorphicForms.targetChart_symm_contMDiffOn, JacobianChallenge.HolomorphicForms.montelForward_contMDiffOn_u_on_patch, JacobianChallenge.HolomorphicForms.montelForward_contMDiffOn_u_on_patch_chartBall}
\leanok
\end{lemma}
\begin{proof}\leanok\end{proof}

\begin{lemma}[Mathlib: atlas member is smooth]
\label{lem:mathlib-contmdiffon-of-mem-maximalatlas}
For an open partial homeomorphism \(e\) in the maximal atlas of a \(C^n\)
manifold, \(e\) is \(\code{ContMDiffOn}\) on \(e.\mathrm{source}\). Green (Mathlib).
\lean{contMDiffOn_of_mem_maximalAtlas}
\leanok
\end{lemma}
\begin{proof}\leanok\end{proof}

\begin{lemma}[Per-patch inverse smoothness of the candidate uniformization]
\label{lem:montel-inverse-per-patch-smoothness}
\uses{lem:montel-forward-per-patch-smoothness, lem:mathlib-contmdiffon-of-mem-maximalatlas, lem:montel-realization-coord-invcoord}
The inverse mirror of the per-patch forward reading — the last local ingredient
underneath the inverse-smoothness leaf
(\(\code{genusZeroMontel\_raw\_global\_patch\_family\_contMDiff\_invMap}\)). If a
Montel patch family represents the inverse branches of a candidate
\(u:X\simeq\code{OnePoint}\,\C\) in local public target charts (\texttt{hinv}) and
each inverse branch is \(\code{ContMDiffOn}\) on its target-chart target
(\texttt{hinv\_smooth}, an honest input — \(\code{invCoord}\) smoothness is not a
patch field, unlike the forward \(\code{coord\_contMDiffOn}\)), then
\(u^{-1}\) is \(\code{ContMDiffOn}\) on each patch's target-chart source. On the
source, \(u^{-1}=\code{invCoord}\circ\code{targetChart}\) (from \texttt{hinv} via
\(\code{targetChart.left\_inv}\)); the forward chart is \(\code{ContMDiffOn}\) as
an atlas member (\code{targetChart\_contMDiffOn}). Selector-free.
\lean{JacobianChallenge.HolomorphicForms.targetChart_contMDiffOn, JacobianChallenge.HolomorphicForms.montelInverse_contMDiffOn_uSymm_on_patch}
\leanok
\end{lemma}
\begin{proof}\leanok\end{proof}

\begin{lemma}[Montel coherent local selector (open)]
\label{lem:montel-coherent-local-selector}
\uses{lem:montel-local-patch-realization, lem:montel-overlap-agreement}
The selected local patches must be coherent with the global selector:
their assigned target charts agree with the family target charts, every global
patch is realized by its local chart-ball data, and local transition
compatibility identifies overlapping coordinates as the same
\(\code{OnePoint}\,\C\) point.
\lean{JacobianChallenge.HolomorphicForms.GenusZeroLocalMontelChartSelector.localPatch_targetChart_eq, JacobianChallenge.HolomorphicForms.GenusZeroCoherentLocalMontelChartSelector, JacobianChallenge.HolomorphicForms.GenusZeroTransitionCoherentLocalMontelChartSelector}
\end{lemma}

\begin{lemma}[Montel overlap agreement kernel]
\label{lem:montel-overlap-agreement}
\uses{lem:montel-overlap-agreement-same-chart, lem:montel-overlap-agreement-cross-chart}
The normalized Montel chart-ball realizations of the finite patch family
agree on patch overlaps.
\lean{JacobianChallenge.HolomorphicForms.genusZeroMontel_normalized_limits_agree_on_overlaps}
\leanok
\end{lemma}

\begin{lemma}[Montel limit uniqueness]
\label{lem:montel-limit-uniqueness}
Two locally-uniform holomorphic limits of the same coordinate function that agree
on a connected overlap must be identical.
\lean{JacobianChallenge.HolomorphicForms.eqOn_of_tendstoLocallyUniformlyOn_same}
\leanok
\end{lemma}

\begin{lemma}[Montel overlap agreement same-chart case (open)]
\label{lem:montel-overlap-agreement-same-chart}
\uses{lem:montel-limit-uniqueness}
When the target charts are the same, the two realizations represent the same
coordinate on the overlap, reducing agreement to the identity theorem.
\lean{JacobianChallenge.HolomorphicForms.genusZeroMontel_overlap_agreement_same_chart}
\leanok
\end{lemma}

\begin{lemma}[Montel overlap agreement cross-chart case]
\label{lem:montel-overlap-agreement-cross-chart}
\uses{lem:montel-overlap-agreement-same-chart, lem:montel-chart-transition-reduction}
When the target charts differ, conjugate one side through the public transition
to reduce to the same-chart case.
\lean{JacobianChallenge.HolomorphicForms.genusZeroMontel_overlap_agreement_cross_chart}
\leanok
\end{lemma}

\begin{lemma}[Montel chart transition reduction]
\label{lem:montel-chart-transition-reduction}
\uses{lem:onepoint-cx-chart-overlap-derivative}
Conjugate through the public transition \(z \mapsto z^{-1}\) on the overlap of chart targets.
\lean{JacobianChallenge.HolomorphicForms.targetChart_inv_eq, JacobianChallenge.HolomorphicForms.targetChart_symm_inv_eq}
\leanok
\end{lemma}

\begin{lemma}[Montel chart-ball source-domain selector (open)]
\label{lem:montel-chartball-source-domain-selector}
\uses{lem:montel-coherent-local-selector}
The selected source domains must be exactly the chart-ball domains of the
attached local Montel coordinates. This is the real
\code{source\_eq\_chartBall\_domain} field of
\code{GenusZeroChartBallDomainLocalMontelSelector}; the missing provider is
the global selection proof producing such a selector.
\lean{JacobianChallenge.HolomorphicForms.GenusZeroChartBallDomainLocalMontelSelector.source_eq_chartBall_domain}
\end{lemma}

\begin{lemma}[Montel two-chart target assignment (open)]
\label{lem:montel-two-chart-target-assignment}
\uses{lem:montel-chartball-source-domain-selector}
The finite selector must reduce to the two public target charts: an identity
patch, an inversion patch, their target-chart identifications, and exhaustion
of all selected indices by these two patches.
\lean{JacobianChallenge.HolomorphicForms.GenusZeroTwoChartMontelAtlasSelector}
\end{lemma}

\begin{lemma}[Mathlib: finite subcover of a compact space from neighborhoods]
\label{lem:mathlib-compactspace-elim-nhds-subcover}
\lean{CompactSpace.elim_nhds_subcover}
On a compact space, from a choice of neighborhood \(U_x\in\mathcal N(x)\) for each
point one can extract a \emph{finite} set of points \(t\) with
\(\bigcup_{x\in t}U_x=\top\). This is the Mathlib compactness primitive behind the
``choose finitely many chart-ball patches'' selection step; the companion
\code{IsCompact.elim\_finite\_subcover} is its open-cover form. Green (Mathlib).
\leanok
\end{lemma}

\begin{lemma}[Montel two-chart source cover obligation (open)]
\label{lem:montel-two-chart-source-cover}
\uses{lem:montel-two-chart-target-assignment, lem:mathlib-compactspace-elim-nhds-subcover}
The identity and inversion local Montel patches must cover the source
manifold \(X\). This is not an accessor lemma: it is the compactness/selection
obligation that proves the chosen chart-ball domains cover \(X\). Each point
has a normalized Montel chart-ball neighborhood, so
\(\code{CompactSpace.elim\_nhds\_subcover}\)
(Lemma~\ref{lem:mathlib-compactspace-elim-nhds-subcover}) extracts a finite
subcover, which the genus-zero normalization then reduces to the two public
(identity / inversion) patches. This node stays open: only the named Mathlib
floor of the selection step is now wired; the genus-zero reduction to exactly two
charts remains an unproven assembly.
\end{lemma}

\begin{lemma}[Montel two-chart target cover (open)]
\label{lem:montel-two-chart-target-cover}
\uses{lem:montel-two-chart-target-assignment}
The identity and inversion public charts must cover
\(\code{OnePoint}\,\C\) on the target side: every target point lies in
\code{identityChart.source} or in \code{inversionChart.source}.
\lean{JacobianChallenge.HolomorphicForms.onePointCx_identity_or_inversionChart_source}
\leanok
\end{lemma}

\begin{lemma}[Montel normalized-limit two-chart cover with realized patches (open)]
\label{lem:montel-normalized-limit-two-chart-cover-with-realized-patches}
\uses{lem:montel-two-chart-source-cover, lem:montel-realized-patch-bundle}
The finite Montel data chooses the normalized identity and inversion local
Montel chart patches, together with realized-patch bundles showing that each
returned global patch is constructed from the corresponding chart-ball limit.
This is the finite-cover frontier before forgetting the packaged
\code{MontelRealizedPatch} data to raw local realization witnesses.
\lean{JacobianChallenge.HolomorphicForms.genusZeroMontel_finite_normalized_chartBall_cover_with_realized_patches}
\end{lemma}

\begin{lemma}[Montel normalized-limit two-chart cover with local realizations (open)]
\label{lem:montel-normalized-limit-two-chart-cover-with-local-realizations}
\uses{lem:montel-normalized-limit-two-chart-cover-with-realized-patches}
The realized-patch finite Montel data forgets each
\code{MontelRealizedPatch} to its raw
\code{GenusZeroLocalMontelPatchRealization} witness while preserving the same
finite family, source cover, target-chart cover, and two public target-chart
assignment. This is the finite-cover frontier before forgetting local
realizations to raw chart balls and source-coordinate maps.
\lean{JacobianChallenge.HolomorphicForms.genusZeroMontel_finite_normalized_chartBall_cover_with_local_realizations}
\end{lemma}

\begin{lemma}[Montel normalized-limit two-chart cover without candidate (open)]
\label{lem:montel-normalized-limit-two-chart-cover-without-candidate}
\uses{lem:montel-normalized-limit-two-chart-cover-with-local-realizations}
The local-realization-backed finite Montel data forgets to the projected
chart-ball packages and source-coordinate maps for the returned global patch
family, but does not yet carry a candidate homeomorphism or prove that those
coordinates represent it globally.
\lean{JacobianChallenge.HolomorphicForms.genusZeroMontel_finite_normalized_chartBall_cover_without_candidate}
\end{lemma}

\begin{lemma}[Montel normalized-limit two-chart cover without coordinate representation (open)]
\label{lem:montel-normalized-limit-two-chart-cover-without-coord-rep}
\uses{lem:montel-normalized-limit-two-chart-cover-without-candidate}
The finite-cover data together with the input homeomorphism
\(_e : X \simeq_t \code{OnePoint}\,\C\) gives the normalized two-chart cover
before coordinate representation. The candidate is reinserted here, keeping the
finite-cover selection node independent of it.
\lean{JacobianChallenge.HolomorphicForms.genusZeroMontel_finite_normalized_chartBall_cover_without_coord_representation}
\end{lemma}

\begin{lemma}[Stage marked topological control data (open)]
\label{lem:stage-marked-topological-control-data}
The Perron engine uses the input homeomorphism
\(_e : X \simeq_t \code{OnePoint}\,\C\) only to choose topological control
data: the two marked source points corresponding to \(0\) and \(\infty\), a
base normalization point away from them, and separated chart neighborhoods
around all three points.  This node is topological bookkeeping for the later
analytic stage construction; it does not assert that \(_e\) is holomorphic.
\end{lemma}

\begin{lemma}[Stage bordered exhaustion domains (open)]
\label{lem:stage-bordered-exhaustion-domains}
\uses{lem:stage-marked-topological-control-data}
There is an increasing sequence of bordered or chart-polygon stage domains
\(\Omega_n\subset X\), adapted to the marked data, whose compact interiors
eventually contain every selected Montel source compactum and whose boundary is
controlled by finitely many source charts.  These domains are the non-circular
places where the analytic stage maps will be constructed before being extended
by filler values outside \(\Omega_n\).
\end{lemma}

\begin{lemma}[Stage cut-domain conjugate control (open)]
\label{lem:stage-cut-domain-conjugate-control}
\uses{lem:stage-bordered-exhaustion-domains}
For each bordered stage domain one chooses cuts giving a cut stage
\(\Omega_n^{\mathrm{cut}}\) on which the harmonic conjugates used by the
Perron engine are single-valued.  The cuts are chosen compatibly with the
exhaustion so every Montel chart-ball compactum is eventually disjoint from the
cuts and contained in the holomorphic part of the stage.
\end{lemma}

\begin{lemma}[Stage dipole boundary control (open)]
\label{lem:stage-dipole-boundary-control}
\uses{lem:stage-marked-topological-control-data, lem:stage-bordered-exhaustion-domains}
The stage boundary data is built from logarithmic dipole models at the two
marked ends, with boundary normalization and compact-subdomain bounds suitable
for Perron/Dirichlet solving.  Existing model-plane and chart-local dipole
lemmas supply only the local singular profiles; the bordered-stage boundary
problem remains new engine work.
\end{lemma}

\begin{lemma}[Stage Dirichlet harmonic solution (open)]
\label{lem:stage-dirichlet-harmonic-exists}
\uses{lem:stage-bordered-exhaustion-domains, lem:stage-dipole-boundary-control}
On each cut stage domain, solve the normalized Dirichlet/Perron problem for
the real harmonic dipole potential with the prescribed boundary control and
marked-end singular behavior.  This is the load-bearing potential-theory
frontier of the engine: Mathlib supplies harmonic predicates and Poisson
formulas for already-harmonic functions, but not this bordered-stage existence
theorem.
\end{lemma}

\begin{lemma}[Stage harmonic conjugate on cuts (open)]
\label{lem:stage-harmonic-conjugate-exists-on-cut}
\uses{lem:stage-cut-domain-conjugate-control, lem:stage-dirichlet-harmonic-exists}
The harmonic solution on each cut stage admits a globally compatible harmonic
conjugate on \(\Omega_n^{\mathrm{cut}}\).  Local harmonic-conjugate and
dipole-conjugate lemmas supply chart-level models; this node carries the
period-killing and compatibility work needed to make the conjugate
single-valued on the chosen cut stage.
\end{lemma}

\begin{lemma}[Stage holomorphic coordinate from harmonic dipole (open)]
\label{lem:stage-holomorphic-coordinate-from-harmonic-dipole}
\uses{lem:stage-dirichlet-harmonic-exists, lem:stage-harmonic-conjugate-exists-on-cut}
Combining the harmonic potential with its conjugate, then exponentiating or
integrating the resulting complex differential, gives a holomorphic
coordinate-like stage map on the cut stage domain.  This is the analytic
production output before target-chart normalization and before any Montel
compactness extraction.
\end{lemma}

\begin{lemma}[Stage map normalization and noncollapse (open)]
\label{lem:stage-map-normalized-noncollapse}
\uses{lem:stage-marked-topological-control-data, lem:stage-holomorphic-coordinate-from-harmonic-dipole}
Normalize the holomorphic stage maps at the fixed marked source data so the
identity-chart local reading has center value \(0\) and derivative \(1\), with
the inversion-chart patch normalized compatibly at the other marked end.  This
exact normalization is the invariant that prevents collapse to a constant in
the later normal-family limit.
\end{lemma}

\begin{lemma}[Stage chart-reading uniform Cauchy bounds (open)]
\label{lem:stage-chart-reading-uniform-cauchy-bound}
\uses{lem:stage-map-normalized-noncollapse, lem:stage-dipole-boundary-control}
The normalized stage maps have uniform bounds on the Cauchy circles of every
eventually contained selected chart ball.  These bounds are the analytic
bridge from Perron stage maps to the green local Montel estimates; they are
proved from boundary control, maximum-principle estimates, and the chosen
normalization.
\end{lemma}

\begin{lemma}[Stage chart readings as ChartBallPowerSeries (open)]
\label{lem:stage-chart-reading-chartBallPowerSeries}
\uses{lem:stage-cut-domain-conjugate-control, lem:stage-holomorphic-coordinate-from-harmonic-dipole, lem:stage-chart-reading-uniform-cauchy-bound}
For each selected Montel patch and all sufficiently large stages, the target
chart reading of the total stage map is packaged on the source chart ball as a
\code{ChartBallPowerSeries}.  This is where the engine enters the existing
chart-ball API: after the eventual containment and holomorphicity hypotheses
are supplied, the local packaging is the standard closed-ball continuous,
open-ball differentiable chart-reading step.
\end{lemma}

\begin{lemma}[Stage common diagonal subsequence (open)]
\label{lem:stage-common-diagonal-subsequence}
\uses{lem:stage-chart-reading-uniform-cauchy-bound, lem:stage-chart-reading-chartBallPowerSeries}
The finitely many stage chart-ball families have a single subsequence that
converges locally uniformly on all selected source chart balls.  The green
finite diagonal extraction and local Montel compactness declarations consume
the chart-ball packages and uniform Cauchy bounds produced by the engine.
\end{lemma}

\begin{lemma}[Stage total filler locally uniform irrelevance (open)]
\label{lem:stage-total-filler-locally-uniform-irrelevance}
\uses{lem:stage-bordered-exhaustion-domains, lem:stage-common-diagonal-subsequence}
The arbitrary filler values used to make each partial stage map a total
function \(F_n:X\to\code{OnePoint}\,\C\) do not affect locally uniform
convergence on any selected patch source, because every compact source
subball is eventually contained in the holomorphic stage domain.  This node is
routine filter and source-local bookkeeping once the exhaustion containment
facts are available.
\end{lemma}

\begin{lemma}[Stage engine common sequence payload (open)]
\label{lem:stage-engine-common-sequence-payload}
\uses{lem:stage-common-diagonal-subsequence, lem:stage-total-filler-locally-uniform-irrelevance, lem:stage-map-normalized-noncollapse}
The Perron engine produces the exact common target sequence payload demanded
by the realized-patch coordinate-representation provider: a total sequence
\(F:\mathbb N\to X\to\code{OnePoint}\,\C\) whose target-chart readings
converge locally uniformly to the realized patch coordinates on every selected
patch source.  The normalization data also supplies the nonconstant local
limits used by the realized Montel patches.
\end{lemma}

\begin{lemma}[Stage engine provider final edge (open)]
\label{lem:stage-engine-provider-final-edge}
\uses{lem:stage-engine-common-sequence-payload, lem:montel-normalized-limit-two-chart-cover-with-realized-patches}
The common sequence payload from the Perron engine is combined with the
realized two-chart Montel cover, overlap agreement, and local coordinate
representation bookkeeping to supply the full realized-patch coordinate
representation provider.  This is the explicit terminal edge from the engine
subtree into the existing Montel frontier.
\end{lemma}

\begin{lemma}[Montel finite-cover coordinate representation with realized patches (open)]
\label{lem:montel-finite-cover-coordinate-representation-with-realized-patches}
\uses{lem:montel-normalized-limit-two-chart-cover-with-realized-patches, lem:montel-coherent-local-selector, lem:stage-engine-provider-final-edge}
The coherent realized-patch finite-cover data constructs the candidate
uniformization existentially, carries the common target sequence payload
demanded by the coherent local selector's overlap-agreement kernel, then proves
that the selected realized patch coordinates represent the candidate in each
public target chart and that the inverse branches represent the candidate
inverse. This is the gluing/coherence frontier separated from the finite-cover
selection step while still exposing the realized-patch bundle data.
\lean{JacobianChallenge.HolomorphicForms.genusZeroMontel_finite_cover_coord_representation_with_realized_patches}
\end{lemma}

\begin{lemma}[Montel finite-cover coordinate representation (open)]
\label{lem:montel-finite-cover-coordinate-representation}
\uses{lem:montel-finite-cover-coordinate-representation-with-realized-patches}
The realized-patch coordinate-representation provider forgets each
\code{MontelRealizedPatch} to its raw
\code{GenusZeroLocalMontelPatchRealization} witness and discards the common
target sequence payload. The public provider preserves the existing downstream
statement: it returns the candidate uniformization, the raw global patch family,
local realization witnesses, and the coordinate/inverse-coordinate
representation facts.
\lean{JacobianChallenge.HolomorphicForms.genusZeroMontel_finite_cover_coord_representation}
\end{lemma}

\begin{lemma}[Montel normalized-limit two-chart cover without source cover (open)]
\label{lem:montel-normalized-limit-two-chart-cover-without-source-cover}
\uses{lem:montel-finite-cover-coordinate-representation}
Projecting chart balls and source-coordinate maps from the local realizations
returned by coordinate representation gives the normalized identity/inversion
chart-ball cover. This provider no longer carries a separate source-cover
obligation; the cover is recovered from the returned family's
\code{patch_cover} field.
\lean{JacobianChallenge.HolomorphicForms.genusZeroMontel_finite_normalized_chartBall_cover_without_source_cover}
\end{lemma}

\begin{lemma}[Montel normalized-limit two-chart cover without target cover (open)]
\label{lem:montel-normalized-limit-two-chart-cover-without-target-cover}
\uses{lem:montel-normalized-limit-two-chart-cover-without-source-cover}
The two chosen public-chart patches retain their normalized chart-ball limit
packages, one for the identity chart and one for the inversion chart, so the
finite cover remembers the local Montel normalization used before global
gluing. The source cover is projected from the returned global patch family's
\code{patch_cover} field. This raw provider omits the target-cover fact, which
is already the green public-chart cover
Lemma~\ref{lem:montel-two-chart-target-cover}.
\lean{JacobianChallenge.HolomorphicForms.genusZeroMontel_finite_normalized_chartBall_cover_without_target_cover}
\end{lemma}

\begin{lemma}[Montel normalized-limit two-chart cover (open)]
\label{lem:montel-normalized-limit-two-chart-cover}
\uses{lem:montel-normalized-limit-two-chart-cover-without-target-cover, lem:montel-two-chart-target-cover}
The finite normalized Montel cover together with the green identity/inversion
target-cover theorem gives the normalized two-chart cover provider used by the
raw patch-family construction.
\lean{JacobianChallenge.HolomorphicForms.GenusZeroNormalizedLimitTwoChartMontelCoverSelector.identity_normalizedLimit, JacobianChallenge.HolomorphicForms.GenusZeroNormalizedLimitTwoChartMontelCoverSelector.inversion_normalizedLimit, JacobianChallenge.HolomorphicForms.genusZeroMontel_finite_normalized_chartBall_cover}
\end{lemma}

\begin{lemma}[Montel global patch-family packaging (open)]
\label{lem:montel-global-patch-family-packaging}
\uses{lem:montel-normalized-limit-two-chart-cover}
The normalized two-chart cover is then forgotten down to the finite
\code{GenusZeroGlobalPatchFamily}: the finite index type, patch family,
source-cover field, and target-chart-cover field. The missing provider is the
construction that packages the normalized-limit cover as this raw family before
the final selector exists.
\end{lemma}

\begin{lemma}[Montel uniformizing patch-family construction (open)]
\label{lem:montel-selector-finite-patch-family-assembly}
\uses{lem:uniformization-genus-zero-homeomorphism, lem:montel-local-patch-realization, lem:montel-coherent-local-selector, lem:montel-chartball-source-domain-selector, lem:montel-two-chart-source-cover, lem:montel-two-chart-target-cover, lem:montel-normalized-limit-two-chart-cover, lem:global-gluing-to-map-exists, lem:global-gluing-chart-expression, lem:global-gluing-local-left-inverse, lem:global-gluing-local-right-inverse, lem:global-gluing-contmdiff-to-map, lem:global-gluing-contmdiff-inv-map}
From a topological homeomorphism \(X\simeq\code{OnePoint}\,\C\) and the
genus-zero hypotheses, construct a finite family of normalized holomorphic
chart-ball patches together with the uniformizing map they determine. This is
the real open frontier: choose the local Montel limits, prove their coherence,
extract the finite source cover by compactness, attach the two public
\(\code{OnePoint}\,\C\) target charts, and feed the resulting compatible data
into the global-gluing substrate. No green Lean declaration currently performs
this construction, so the node intentionally has no \code{\textbackslash lean}
tag.
\lean{JacobianChallenge.HolomorphicForms.genusZeroMontel_raw_global_patch_family}
\end{lemma}

\begin{lemma}[Montel constructed patch family (open)]
\label{lem:montel-selector-patch-family-extraction}
\uses{lem:montel-selector-finite-patch-family-assembly}
The construction above must expose its finite patch family: normalized source
patches, their public \(\code{OnePoint}\,\C\) target charts, a source cover, and
a target-chart cover. This is a property of the constructed data, not a
projection from a pre-existing selector structure.
\end{lemma}

\begin{lemma}[Montel constructed coordinate representation (open)]
\label{lem:montel-selector-coord-represents-uniformization}
\uses{lem:montel-selector-finite-patch-family-assembly, lem:montel-selector-patch-family-extraction}
For the constructed patch family, each selected local coordinate must represent
the constructed global uniformization in its assigned target chart:
\code{targetChart.symm (coord x) = uniformization x}. This property depends on
the construction; it is not an input to it.
\lean{JacobianChallenge.HolomorphicForms.genusZeroMontel_raw_global_patch_family}
\end{lemma}

\begin{lemma}[Montel constructed inverse-coordinate representation (open)]
\label{lem:montel-selector-invcoord-represents-uniformization}
\uses{lem:montel-selector-finite-patch-family-assembly, lem:montel-selector-patch-family-extraction}
For the constructed patch family, each selected inverse coordinate branch must
represent the inverse of the constructed global uniformization on the assigned
target chart:
\code{invCoord z = uniformization.symm (targetChart.symm z)}. This property
depends on the construction and on the inverse branches supplied by the local
Montel chart data.
\lean{JacobianChallenge.HolomorphicForms.genusZeroMontel_raw_global_patch_family}
\end{lemma}

\begin{lemma}[Montel constructed forward smoothness (open)]
\label{lem:montel-selector-uniformization-contmdiff}
\uses{lem:montel-selector-finite-patch-family-assembly, lem:montel-selector-coord-represents-uniformization, lem:global-gluing-contmdiff-to-map}
The constructed global uniformization represented by the selected local
coordinates must be \(\code{ContMDiff}\). This property follows only after the
construction has produced coherent local chart expressions and the global
gluing smoothness theorem applies.
\lean{JacobianChallenge.HolomorphicForms.genusZeroMontel_raw_global_patch_family_contMDiff_toMap}
\end{lemma}

\begin{lemma}[Montel constructed inverse smoothness (open)]
\label{lem:montel-selector-inverse-uniformization-contmdiff}
\uses{lem:montel-selector-finite-patch-family-assembly, lem:montel-selector-invcoord-represents-uniformization, lem:global-gluing-contmdiff-inv-map}
The inverse of the constructed global uniformization represented by the selected
inverse coordinate branches must be \(\code{ContMDiff}\). This property depends
on the constructed inverse branches and the global gluing smoothness theorem for
the inverse map.
\lean{JacobianChallenge.HolomorphicForms.genusZeroMontel_raw_global_patch_family_contMDiff_invMap}
\end{lemma}

\begin{lemma}[Normalized Montel patch selector from a sphere homeomorphism (open)]
\label{lem:montel-selector-of-homeomorph-onepoint}
\uses{lem:montel-selector-finite-patch-family-assembly, lem:montel-selector-patch-family-extraction, lem:montel-selector-coord-represents-uniformization, lem:montel-selector-invcoord-represents-uniformization, lem:montel-selector-uniformization-contmdiff, lem:montel-selector-inverse-uniformization-contmdiff}
From the constructed uniformizing patch family and its property theorems,
package the current Lean-facing
\code{GenusZeroNormalizedMontelPatchSelector}. This node no longer uses the
selector's own field projections as prerequisites: the construction comes first,
and the representation and smoothness facts are downstream properties of that
construction. It is tracked by the current provider
\code{genusZeroNormalizedMontelPatchSelector\_of\_homeomorph\_onePoint}; it
stays open because that provider is the remaining uniformization frontier.
\lean{JacobianChallenge.HolomorphicForms.genusZeroNormalizedMontelPatchSelector_of_homeomorph_onePoint}
\end{lemma}

\begin{lemma}[Global gluing data from the Montel selector (open)]
\label{lem:genus-zero-gluing-data-from-selector}
\uses{lem:montel-selector-of-homeomorph-onepoint, lem:global-gluing-to-map-exists, lem:global-gluing-chart-expression, lem:global-gluing-local-left-inverse, lem:global-gluing-local-right-inverse, lem:global-gluing-contmdiff-to-map, lem:global-gluing-contmdiff-inv-map}
The intermediate Lean step actually called by the \#232 root: from a topological
\(X\simeq\code{OnePoint}\,\C\), produce a completed \code{GenusZeroGlobalGluingData}.
Its body is \emph{direct}-sorry-free --- it discharges every gluing-data field
(target membership, chart expression, overlap compatibility, the two local
inverse laws, both \code{ContMDiff} fields) from the selector's fields and the
green global-gluing lemmas above. It nonetheless stays \emph{open} (it is left
deliberately un-green): it obtains the selector via
\code{genusZeroNormalizedMontelPatchSelector\_of\_homeomorph\_onePoint}, so it is
transitively \code{sorryAx}-tainted through
Lemma~\ref{lem:montel-selector-of-homeomorph-onepoint}. This node makes the real
selector\(\to\)gluing-data call edge graph-visible, instead of the root appearing
to depend on the selector and the green lemmas directly.
\lean{JacobianChallenge.HolomorphicForms.genusZeroGlobalGluingData_of_homeomorph_onePoint}
\end{lemma}

\begin{theorem}[Biholomorphism from sphere homeomorphism (genus-zero uniformization)]
\label{thm:uniformization-genus-zero-biholomorphism}
\uses{lem:genus-zero-gluing-data-from-selector}
A compact connected complex 1-manifold homeomorphic to \(S^2\) admits a
\emph{biholomorphism} to \(\code{OnePoint}\,\C\): a homeomorphism
\(f:X\simeq\code{OnePoint}\,\C\) with both \(f\) and \(f^{-1}\) complex-smooth
(\code{ContMDiff}).
Tracking: \ghissue{232}.
\lean{JacobianChallenge.HolomorphicForms.exists_biholomorph_onePoint_of_genus_zero}
\notready
\end{theorem}
\begin{proof}
This is the genus-zero uniformization theorem. Classically it follows from the
Riemann mapping theorem with a removable-singularities argument, or from
Riemann--Roch (genus \(0\Rightarrow\) a degree-one meromorphic map to
\(\code{OnePoint}\,\C\)). \emph{The Lean development takes neither route}: the
Riemann--Roch section chain (\code{genusZero\_pointRRSection\_*}) is circular
here, feeding back through the genus-zero single-pole data it is meant to
produce. Instead the biholomorphism is built from a self-contained local
analytic uniformization substrate --- local chart-ball power series, an
open-mapping / inverse-function normalization of local coordinates, a normal
family (Montel / Arzel\`a--Ascoli) extraction of compatible normalized patches,
and a gluing through the two-chart \(\code{OnePoint}\,\C\) atlas that is
\code{ContMDiff} in both directions. The green global-gluing lemmas above show
how a completed \code{GenusZeroNormalizedMontelPatchSelector} produces the
forward map, inverse branches, local inverse laws, and \code{ContMDiff}
regularity. The remaining orange frontier is the construction of such a
selector from the normal-family substrate, packaged by this theorem. The
simple-pole / degree-one meromorphic data of the classical proof is then
\emph{obtained from} this biholomorphism, not used to construct it.
\end{proof}

\subsection{Path B: Riemann--Roch single-pole map decomposition}
\label{sec:path-b-rr-single-pole-map}

\begin{definition}[Germ-space carrier for Riemann--Roch sections]
\label{def:rr-germ-space-carrier}
The sound module-facing carrier for \(L(D)\) is the submodule of meromorphic
germ families generated by divisor-compatible bounded sections. Arithmetic is
performed in local \(\code{Filter.Germ}\)s, so pole cancellation is represented
in the carrier rather than by pointwise arithmetic on \(\code{OnePoint}\,\C\).
\lean{JacobianChallenge.HolomorphicForms.RiemannRochGermSpace}
\leanok
\end{definition}

\begin{definition}[Genus-zero point RR germ-space dimension input]
\label{def:rr-germ-space-dimension-input}
\uses{def:rr-germ-space-carrier}
The open hard input for Path B is the interface asserting
\[
  \code{finrank}\,\C\,
    (\code{RiemannRochGermSpace}\,X\,(\code{Divisor.point}\,P)) = 2,
\]
together with the residual \(K-[P]\) vanishing term. Inhabiting this interface
is the actual Riemann--Roch dimension theorem for the sound germ-space carrier;
this node is intentionally left without \(\code{\textbackslash leanok}\).
\lean{JacobianChallenge.HolomorphicForms.GenusZeroPointRiemannRochGermSpaceDimensionInput}
\end{definition}

\begin{lemma}[Genus-zero point RR germ-space dimension provider]
\label{lem:rr-germ-space-dimension-input-provider}
\uses{def:rr-germ-space-dimension-input}
The named open provider for the Path B dimension layer inhabits the point
Riemann--Roch germ-space dimension-input interface. This is the remaining
Riemann--Roch dimension theorem for the sound germ-space carrier, and it is
intentionally left open.
\lean{JacobianChallenge.HolomorphicForms.genusZero_pointRiemannRochGermSpaceDimensionInput}
\end{lemma}

\begin{lemma}[Constant germ line in the RR germ space]
\label{lem:rr-germ-space-constant-line}
\uses{def:rr-germ-space-carrier}
Constants form a one-dimensional submodule of meromorphic germ families, and
for every effective divisor \(D\) this constant line lies in
\(\code{RiemannRochGermSpace}\,X\,D\). This is the constant-line input used by
the linear-algebra lemma comparing a two-dimensional ambient space with its
one-dimensional subspace of constants.
\lean{JacobianChallenge.HolomorphicForms.constantGermFamilyLine}
\lean{JacobianChallenge.HolomorphicForms.constantGermFamilyLine_le_RiemannRochGermSpace}
\lean{JacobianChallenge.HolomorphicForms.constantGermFamilyLine_finrank}
\leanok
\end{lemma}
\begin{proof}
\leanok
\end{proof}

\begin{lemma}[RR germ outside constants from the finrank input]
\label{lem:rr-germ-space-outside-constants}
\uses{def:rr-germ-space-dimension-input, lem:rr-germ-space-constant-line}
Given the point RR germ-space finrank input, the one-dimensional constant line
is a proper submodule of
\(\code{RiemannRochGermSpace}\,X\,(\code{Divisor.point}\,P)\). Hence there is
an RR germ-space element whose underlying meromorphic germ family is not
constant.
\lean{JacobianChallenge.HolomorphicForms.genusZero_pointRiemannRochGermSpace_exists_outside_constants}
\leanok
\end{lemma}
\begin{proof}
\leanok
\end{proof}

\begin{lemma}[RR germ-to-map representation bridge]
\label{lem:rr-germ-to-map-bridge}
\uses{lem:rr-germ-space-outside-constants}
The remaining representation bridge converts an outside-constant element of the
point RR germ-space carrier into an honest
\(\code{GenusZeroPointRiemannRochElement}\), i.e. a nonconstant meromorphic
map to the sphere in \(L([P])\). This is the current narrow open frontier:
the abstract span element must be represented by divisor-compatible
meromorphic data.
\lean{JacobianChallenge.HolomorphicForms.genusZero_pointRiemannRochElement_of_germSpace_outside_constants}
\end{lemma}

\begin{lemma}[RR element from the germ-space finrank input]
\label{lem:rr-element-of-germ-space-dimension-input}
\uses{def:rr-germ-space-dimension-input, lem:rr-germ-to-map-bridge}
Assuming the point RR germ-space finrank input and the representation bridge,
there is a nonconstant Riemann--Roch element in \(L([P])\). This node remains
open pending the germ-to-map bridge, since the Lean theorem depends
transitively on that open provider.
\lean{JacobianChallenge.HolomorphicForms.genusZero_pointRiemannRochElement_of_germSpaceDimensionInput}
\end{lemma}

\begin{lemma}[Hard core: dimension of \(L([P])\) in genus zero (open)]
\label{lem:rr-dim-L-point-eq-two-core}
\uses{def:rr-germ-space-carrier, lem:rr-germ-space-dimension-input-provider}
The irreducible Riemann--Roch input for the Path B fixed-pole route is the
dimension count
\[
  \dim_\C L([P]) = 2
\]
for every \(P:X\) when \(\mathrm{analyticGenus}_{\C}(X)=0\). A precise future
Lean-facing theorem should expose the actual finite-dimensional
Riemann--Roch space, not merely an existential proxy:
\[
\begin{gathered}
  (X : \mathrm{Type}) \to [\cdots] \to
  (P : X) \to (h : \code{analyticGenus}\,\C\,X = 0) \to\\
  \code{finrank}\,\C\,(\code{RiemannRochSpace}\,X\,(\code{Divisor.point}\,P)) = 2.
\end{gathered}
\]
The current Lean declaration below is only the structural arithmetic placeholder
for the Riemann--Roch difference term, now open pending the RR dimension input;
it does \emph{not} constitute a completed RR vector-space API or the actual
dimension theorem. This is the single hard core that must be built from the
project's divisor, meromorphic-function, and \(\code{analyticGenus}\) /
holomorphic-one-form finiteness machinery.
\lean{JacobianChallenge.HolomorphicForms.genusZero_riemannRoch_difference_eq_two}
\end{lemma}

\begin{lemma}[Vanishing of the residual \(K-[P]\) term (supporting)]
\label{lem:rr-K-minus-point-vanishes}
\uses{lem:rr-germ-space-dimension-input-provider}
The complementary Riemann--Roch term has dimension zero in genus zero:
\[
  \dim_\C L(K-[P]) = 0.
\]
In the current Lean scaffold this is represented by the structural placeholder
\[
\code{genusZero\_riemannRoch\_K\_minus\_point\_dim\_zero}
  : \exists \ell_{K-P}:\mathbb{N},\ \ell_{K-P}=0.
\]
An honest proof should reduce this to negative-degree vanishing for
meromorphic sections of \(K-[P]\), using \(g_{\mathrm{an}}(X)=0\) to identify
\(\deg K=-2\). The current Lean declaration is open pending the RR dimension
input.
\lean{JacobianChallenge.HolomorphicForms.genusZero_riemannRoch_K_minus_point_dim_zero}
\end{lemma}

\begin{lemma}[Dimension \(>1\) gives a nonconstant element of \(L([P])\)]
\label{lem:rr-dim-ge-two-nonconstant}
\uses{lem:rr-dim-L-point-eq-two-core, lem:rr-K-minus-point-vanishes, lem:rr-element-of-germ-space-dimension-input}
Once the real \(\dim_\C L([P])=2\) theorem is available, pure linear algebra
should produce a nonconstant \(f\in L([P])\): constants give a one-dimensional
subspace, so a two-dimensional ambient space contains an element outside the
constant line. The intended Lean shape is the existing theorem
\[
\begin{gathered}
\code{riemannRochSpace\_dim\_ge\_two\_implies\_nonconstant\_meromorphic}\\
  (X) (P) (h : \code{analyticGenus}\,\C\,X = 0) :
  \exists f : \code{MeromorphicMapToSphere}\,X,\ 
    f.\code{Nonconstant} \wedge
    f.\code{MemRiemannRochSpace}\,(\code{Divisor.point}\,P).
\end{gathered}
\]
The Lean declaration now routes through the point RR germ-space dimension
provider and the germ-to-map representation bridge:
\code{genusZero\_pointRiemannRochGermSpaceDimensionInput} supplies the
dimension input and
\code{genusZero\_pointRiemannRochElement\_of\_germSpaceDimensionInput}
packages the resulting element. The node remains open transitively on those
providers, but no longer depends on the fixed-pole analytic RR provider.
\lean{JacobianChallenge.HolomorphicForms.riemannRochSpace_dim_ge_two_implies_nonconstant_meromorphic}
\end{lemma}

\begin{lemma}[A nonconstant element of \(L([P])\) has pole divisor \([P]\)]
\label{lem:rr-nonconstant-L-point-has-exact-pole}
\uses{lem:rr-dim-ge-two-nonconstant}
If \(f\in L([P])\) is nonconstant on a compact connected genus-zero surface, the
Riemann--Roch membership condition gives \(\mathrm{poles}(f)\le [P]\). The pole
divisor is effective and cannot be zero, since a holomorphic meromorphic map on
a compact connected surface is constant. Hence \(\mathrm{poles}(f)=[P]\).
\lean{JacobianChallenge.HolomorphicForms.genusZero_poleDivisor_eq_point_of_nonconstant_mem_L_point}
\leanok
\end{lemma}
\begin{proof}\leanok\end{proof}

\begin{lemma}[A nonconstant element of \(L([P])\) has analytic order one at \(P\)]
\label{lem:rr-nonconstant-L-point-order-one}
\uses{lem:rr-nonconstant-L-point-has-exact-pole}
If \(f\in L([P])\) is nonconstant and carries honest analytic data, the
chart-local analytic order of \(f\) at \(P\) is exactly \(1\): the pole divisor
is pinned to \([P]\) by
Lemma~\ref{lem:rr-nonconstant-L-point-has-exact-pole}, and at a simple pole the
analytic-data record reads off order one in the inversion chart.
\lean{JacobianChallenge.HolomorphicForms.genusZero_mapAnalyticOrderAt_eq_one_of_nonconstant_mem_L_point}
\leanok
\end{lemma}
\begin{proof}\leanok\end{proof}

\begin{lemma}[A nonconstant element of \(L([P])\) is a branched cover of degree one]
\label{lem:rr-nonconstant-L-point-branched-cover-datum}
\uses{lem:rr-nonconstant-L-point-order-one}
A nonconstant \(f\in L([P])\) with analytic data packages as a
\code{BranchedCoverDataOfPoleDegree}: a branched cover of \(\CPone\) whose
branched degree over \(\infty\) equals
\(\deg(\mathrm{poles}(f))=\deg([P])=1\). The simple-pole constructor consumes
the exact pole divisor and the order-one fact of the preceding nodes.
\lean{JacobianChallenge.HolomorphicForms.genusZero_branchedCoverDataOfPoleDegree_of_nonconstant_mem_L_point}
\leanok
\end{lemma}
\begin{proof}\leanok\end{proof}

\begin{lemma}[Degree-one bijectivity package for a nonconstant element of \(L([P])\)]
\label{lem:rr-nonconstant-L-point-degree-one-package}
\uses{lem:rr-nonconstant-L-point-branched-cover-datum, thm:degree-one-bijectivity-meromorphic}
A nonconstant \(f\in L([P])\) with analytic data yields the full
\code{MeromorphicDegreeOneData} record: \(f\) is continuous, \emph{bijective}
onto \(\CPone\), and its pole divisor has degree one. Continuity comes from the
analytic data (so, unlike the generic assembly, no separate pole-modulus
divergence datum is needed), the degree from \(\deg([P])=1\), and bijectivity
from Theorem~\ref{thm:degree-one-bijectivity-meromorphic} applied to the
branched-cover datum of the preceding node.
\lean{JacobianChallenge.HolomorphicForms.genusZero_meromorphicDegreeOneData_of_nonconstant_mem_L_point}
\leanok
\end{lemma}
\begin{proof}\leanok\end{proof}

\begin{lemma}[Mathlib: compact-to-Hausdorff continuous bijection is a homeomorphism]
\label{lem:mathlib-homeo-of-equiv-compact-to-t2}
\lean{Continuous.homeoOfEquivCompactToT2}
A continuous bijection from a compact space to a Hausdorff space is a closed
map, hence its set-theoretic inverse is continuous and the bijection promotes
to a homeomorphism. This is the Mathlib primitive behind the topological
uniformization payoff of the Path-B chain. Green (Mathlib).
\leanok
\end{lemma}

\begin{lemma}[A nonconstant element of \(L([P])\) realizes the sphere homeomorphism]
\label{lem:rr-nonconstant-L-point-homeomorph-sphere}
\uses{lem:rr-nonconstant-L-point-degree-one-package, lem:mathlib-homeo-of-equiv-compact-to-t2}
A nonconstant \(f\in L([P])\) with analytic data realizes a homeomorphism
\(e : X\simeq\code{OnePoint}\,\C\) whose underlying map \emph{is}
\(f.\code{toMap}\). The carried identification \(e = f.\code{toMap}\) (as
functions) lets downstream consumers transport forward holomorphy of \(f\)
along \(e\) toward the genus-zero biholomorphism (\#232) and supplies the
topological \(X\simeq\code{OnePoint}\,\C\) input of \#233.
\lean{JacobianChallenge.HolomorphicForms.genusZero_homeomorph_onePoint_of_nonconstant_mem_L_point}
\leanok
\end{lemma}
\begin{proof}\leanok
The degree-one bijectivity package gives a continuous bijection
\(f : X\to\code{OnePoint}\,\C\) from a compact space to a Hausdorff space;
Lemma~\ref{lem:mathlib-homeo-of-equiv-compact-to-t2} promotes it to a
homeomorphism with the same underlying map.
\end{proof}

\begin{lemma}[A nonconstant element of \(L([P])\) has effective zero divisor]
\label{lem:rr-nonconstant-L-point-zero-divisor-effective}
\uses{lem:rr-nonconstant-L-point-has-exact-pole}
A nonconstant \(f\in L([P])\) has \(\mathrm{zeros}(f)\ge 0\): at \(P\) the
exact pole divisor \([P]\) forces \(\mathrm{zeros}(f)(P)=0\) by
zero/pole-support disjointness, and away from \(P\) the Riemann--Roch
membership inequality reduces to \(\mathrm{zeros}(f)(Q)\ge 0\) directly. This
is one of the four demanded fields of the open effective granular provider
(Lemma~\ref{lem:rr-effective-granular-provider}), derived from the honest
hypotheses alone.
\lean{JacobianChallenge.HolomorphicForms.genusZero_zeroDivisor_effective_of_nonconstant_mem_L_point}
\leanok
\end{lemma}
\begin{proof}\leanok\end{proof}

\begin{lemma}[A nonconstant element of \(L([P])\) lands in the germ-space carrier]
\label{lem:rr-nonconstant-L-point-germ-space-member}
\uses{lem:rr-nonconstant-L-point-zero-divisor-effective}
A nonconstant \(f\in L([P])\) with analytic data packages as a
divisor-compatible Riemann--Roch bounded section
\(s : \code{RiemannRochBoundedSection}\,X\,[P]\) whose underlying function
\emph{is} \(f\) (the carried identification \(s = f.\code{toMap}\) as
functions). Hence its germs are a member of the sound germ-space module
carrier \(\code{RiemannRochGermSpace}\,X\,[P]\) — the span whose
\(\dim_\C = 2\) computation is the Path-B dimension core. This is the seam
where the consumer chain of the preceding nodes and the germ-space carrier
lane meet on one object.
\lean{JacobianChallenge.HolomorphicForms.genusZero_riemannRochBoundedSection_of_nonconstant_mem_L_point}
\leanok
\end{lemma}
\begin{proof}\leanok
The effectivity lemma supplies the zero-divisor field obligation; the
analytic-data lift supplies the underlying meromorphic function and its local
germ family; the divisor fields are copied, with the recorded order function
given by the principal-divisor coefficient, whose compatibility is exactly
zero/pole disjointness plus the two effectivity facts. Riemann--Roch
membership transfers verbatim since the principal divisors agree.
\end{proof}

\begin{lemma}[Germ families outside the constant line have nonconstant functions]
\label{lem:rr-germs-outside-constants-nonconstant}
\uses{lem:rr-nonconstant-L-point-germ-space-member}
A divisor-compatible meromorphic function whose germ family lies outside the
constant germ-family line is nonconstant as a function to
\(\code{OnePoint}\,\C\). Pointwise-constant functions — including the
everywhere-\(\infty\) case, whose finite \(\code{getD}\,0\) lift is the
constant \(0\) — all have constant germ families inside the line, so the
contrapositive transfers the linear-algebra hypothesis
\(F\notin\code{constantGermFamilyLine}\) of the germ-space extraction into
the map-level \code{Nonconstant} field. This is the seam interface consumed
by the open germ-to-map bridge
(Lemma~\ref{lem:rr-germ-to-map-bridge}) when it represents an
outside-constants span element by an honest carrier function.
\lean{JacobianChallenge.HolomorphicForms.MeromorphicFunctionWithDivisors.nonconstant_of_germs_notMem_constantGermFamilyLine}
\leanok
\end{lemma}
\begin{proof}\leanok
A pointwise-constant function has finite lift constantly \(c.\code{getD}\,0\),
hence constant germ at every point, hence germ family
\(\code{constantGermFamilyLinearMap}\,(c.\code{getD}\,0)\) inside the line;
contrapose.
\end{proof}

\begin{lemma}[A carrier representative yields the Riemann--Roch element with analytic data]
\label{lem:rr-carrier-representative-element}
\uses{lem:rr-nonconstant-L-point-germ-space-member, lem:rr-germs-outside-constants-nonconstant}
A divisor-compatible carrier representative \(g\) in \(L([P])\) whose germ
family lies outside the constant line yields a
\code{GenusZeroPointRiemannRochElement} \emph{with} its \code{AnalyticData},
for any sphere map \(f\) linked to \(g\) by the function identification
\(f.\code{toMap} = g.\code{toFunction}.\code{toFun}\) and the divisor-copy
facts, given the order-soundness side condition (simple pole \(\Rightarrow\)
chart-local order one, now backed by the carrier's order--\code{orderAt}
compatibility field) as a hypothesis. Membership
transfers since the principal divisors agree; nonconstancy is the germ-seam
transfer (Lemma~\ref{lem:rr-germs-outside-constants-nonconstant}); analytic
data is carried by the underlying meromorphic function type except for the
order-one field, which is the hypothesis. The statement-level hypotheses here
are now aligned with the carrier order-soundness field; once the open
germ-to-map bridge (Lemma~\ref{lem:rr-germ-to-map-bridge}) produces the
representative, this packaging lemma supplies the analytic-data payload.
\lean{JacobianChallenge.HolomorphicForms.genusZeroPointRiemannRochElement_with_analyticData_of_carrier_representative}
\leanok
\end{lemma}
\begin{proof}\leanok
Assemble the element from the linked map: membership by the principal-divisor
copy, nonconstancy by the germ-seam transfer along the function
identification, analytic data from \code{toFun\_continuous} and
\code{isMeromorphic} plus the order-soundness hypothesis transported along
the pole-divisor copy.
\end{proof}

\begin{lemma}[Span elements are normalized finite sums of bounded-section generators]
\label{lem:rr-span-expression-extraction}
\uses{lem:rr-nonconstant-L-point-germ-space-member}
Every element of the Riemann--Roch germ space is a finite \(\C\)-linear
combination of germ families of bounded-section generators, and the
expression can be normalized so all coefficients are nonzero; nonzero scalar
multiplication of a generator multiplies its germ family by the scalar
(definitionally). This is the routine linear-algebra packaging of Steps 1--2
of the germ-to-map bridge plan (consumed by the open
Lemma~\ref{lem:rr-germ-to-map-bridge}): the bridge's remaining analytic
content, the pairwise divisor-compatible addition of bounded sections
(Step~3), is deliberately NOT covered by this node.
\lean{JacobianChallenge.HolomorphicForms.RiemannRochGermSpace.exists_finset_sum_generators}
\leanok
\end{lemma}
\begin{proof}\leanok
Unfold the span membership via \code{Submodule.mem\_span\_set'}, choose a
bounded-section preimage for each range member, prune zero coefficients
(their summands vanish), and re-enumerate the surviving finite index set.
\end{proof}

\begin{lemma}[Bounded sections package as sphere maps]
\label{lem:rr-bounded-section-to-sphere-map}
\uses{lem:rr-carrier-representative-element, lem:rr-span-expression-extraction}
A divisor-compatible Riemann--Roch bounded section packages as a meromorphic
map to the sphere with its divisor data and \(L([P])\) membership
transferred, and a bounded section of \(L([P])\) whose germ family lies
outside the constant line yields a \code{GenusZeroPointRiemannRochElement}
— with its \code{AnalyticData} given the order-soundness side condition.
This instantiates the statement-level carrier pipeline with the
carrier-to-sphere wrapper, every linking hypothesis discharged
definitionally; it is the Step-5--7 packaging of the germ-to-map bridge plan
(Lemma~\ref{lem:rr-germ-to-map-bridge}), whose remaining content is thereby
exactly its Steps 1--4: the span expression
(Lemma~\ref{lem:rr-span-expression-extraction}) and the pairwise
divisor-compatible addition fold. (Conditional on the carrier
order-soundness gap — see the CARRIER-GATED
Lemma~\ref{lem:germ-order-bridge}: the wrapper's pole-value field currently
routes through that gap, so the declarations here inherit its open status;
when the carrier migration lands, this node greens automatically with zero
edits.)
\lean{JacobianChallenge.HolomorphicForms.genusZeroPointRiemannRochElement_with_analyticData_of_boundedSection}
\end{lemma}

\begin{lemma}[A nonconstant element of \(L([P])\) satisfies the granular field matrix]
\label{lem:rr-nonconstant-L-point-granular-fields}
\uses{lem:rr-nonconstant-L-point-zero-divisor-effective, lem:rr-nonconstant-L-point-order-one}
A nonconstant \(f\in L([P])\) with analytic data satisfies the exact
four-field matrix demanded by the open effective granular provider
(Lemma~\ref{lem:rr-effective-granular-provider}): effective zero divisor,
pole divisor exactly \([P]\), pointwise meromorphic canonical finite lift,
and order-one extension at the pole. The accompanying existential reduction
\code{genusZero\_fixedPole\_rr\_effectiveGranular\_of\_element} states the
provider's literal goal conditional on an element plus its analytic data, so
the remaining Path-B gap is precisely the extraction frontier: produce a
\code{GenusZeroPointRiemannRochElement} together with its
\code{AnalyticData}.
\lean{JacobianChallenge.HolomorphicForms.genusZero_effectiveGranular_fields_of_nonconstant_mem_L_point}
\leanok
\end{lemma}
\begin{proof}\leanok
Conjunction of the zero-divisor effectivity and exact-pole-divisor leaves,
the \code{meromorphic\_getD} field of the analytic data, and the order-one
bridge.
\end{proof}

\begin{lemma}[Path B fixed-pole effective granular provider (open)]
\label{lem:rr-effective-granular-provider}
\uses{lem:rr-nonconstant-L-point-has-exact-pole}
The direct Path B frontier packages the Riemann--Roch construction as a
\(\code{MeromorphicMapToSphere}\) with the exact simple-pole divisor and the
granular analytic fields needed downstream:
\[
\begin{gathered}
\exists f : \code{MeromorphicMapToSphere}\,X,\ 
  \code{Divisor.Effective}\,f.\code{zeroDivisor} \wedge
  f.\code{poles} = \code{Divisor.point}\,P \wedge\\
  (\forall p:X,\ \code{MeromorphicAtX}\,(\lambda q,\ (f.\code{toMap}\,q).\code{getD}\,0)\,p)
  \wedge
  \code{mapAnalyticOrderAt}\,f.\code{toMap}\,P = 1.
\end{gathered}
\]
This provider should be built from the nonconstant \(L([P])\) element and the
local Laurent/order analysis of that RR section. It must not use the already
packaged single-pole data, the biholomorphism route, or any field projection
from the structure it is trying to construct.
\lean{JacobianChallenge.HolomorphicForms.genusZero_fixedPole_rr_effectiveGranular_provider}
\end{lemma}

\begin{lemma}[Effective granular provider packages the analytic RR witness]
\label{lem:rr-effective-granular-packages-analytic-data}
\uses{lem:rr-effective-granular-provider}
From the effective granular provider, the existing sorry-free assembly derives
Riemann--Roch membership, nonconstancy, continuity of the map to
\(\code{OnePoint}\,\C\), the \(\code{AnalyticData}\) record, and pole-modulus
divergence. The current Lean chain is
\[
\code{genusZero\_fixedPole\_rr\_granularAnalytic\_provider}
\to
\code{genusZero\_fixedPole\_rr\_granularMap\_provider}
\to
\code{genusZero\_fixedPole\_rr\_analyticMap\_provider}
\to
\code{genusZero\_fixedPole\_rr\_analyticData\_provider}.
\]
Thus the only open construction beneath this packaging chain is
Lemma~\ref{lem:rr-effective-granular-provider}.
\lean{JacobianChallenge.HolomorphicForms.genusZero_fixedPole_rr_analyticData_provider}
\end{lemma}

\begin{lemma}[Analytic genus zero gives a sphere homeomorphism]
\label{lem:analytic-genus-zero-homeomorphic-sphere}
\uses{lem:rr-effective-granular-packages-analytic-data, lem:meromorphic-to-branched-cover-data, thm:degree-one-bijective}
If a compact connected Riemann surface has analytic genus zero, then it is
homeomorphic to the standard \(2\)-sphere.
\lean{JacobianChallenge.HolomorphicForms.homeomorphic_sphere_of_analyticGenus_eq_zero}
\leanok
\end{lemma}
\begin{proof}\leanok
This is the topological classification route: construct a genus-zero
single-pole meromorphic map, use the meromorphic-to-branched-cover bridge and
degree-one bijectivity to identify \(X\) homeomorphically with
\(\code{OnePoint}\,\C\), then compose with the fixed homeomorphism from
\(\code{OnePoint}\,\C\) to \(S^2\).
\end{proof}

\begin{lemma}[Analytic genus zero gives a sphere biholomorphism (open)]
\label{lem:exists-biholomorph-analyticgenus-zero}
\uses{lem:analytic-genus-zero-homeomorphic-sphere, thm:uniformization-genus-zero-biholomorphism}
A compact connected Riemann surface with \(g_{\mathrm{an}}(X)=0\) admits a
biholomorphism \(X\simeq\code{OnePoint}\,\C\). This is the thin packaging
layer that first gets the topological sphere homeomorphism from
Lemma~\ref{lem:analytic-genus-zero-homeomorphic-sphere}, then feeds it to
Theorem~\ref{thm:uniformization-genus-zero-biholomorphism}. Open obligation
(\lstinline{sorry}) in the Lean development.
Tracking: \ghissue{233}.
\lean{JacobianChallenge.HolomorphicForms.exists_biholomorph_onePoint_of_analyticGenus_zero}
\end{lemma}

\begin{definition}[Complex simple-pole principal part]
\label{def:complex-simple-pole-principal-part}
For a complex function \(F:X\to\C\) and a point \(P\), the predicate
\(\code{HasComplexSimplePolePrincipalPart}\,F\,P\) bundles the four analytic
facts needed for a simple-pole map to the sphere: meromorphicity everywhere,
continuity of the \(\code{OnePoint}\)-extension, chart-local order \(1\) at
\(P\), and norm divergence along the punctured neighbourhood of \(P\).
\lean{JacobianChallenge.HolomorphicForms.HasComplexSimplePolePrincipalPart}
\leanok
\end{definition}

\begin{definition}[Simple-pole sphere data from a complex principal part]
\label{def:simple-pole-to-sphere-data-of-complex-principal-part}
\uses{def:complex-simple-pole-principal-part}
A complex simple-pole principal part \(F\) at \(P\) canonically packages as
\(\code{SimplePoleToSphereData}\) by using the \(\code{OnePoint}\)-extension of
\(F\) as the map to the sphere.
\lean{JacobianChallenge.HolomorphicForms.SimplePoleToSphereData.of_complexPrincipalPart}
\leanok
\end{definition}

\begin{lemma}[Simple-pole principal part from a sphere biholomorphism (open)]
\label{lem:complex-simple-pole-principal-part-biholomorph}
\uses{def:complex-simple-pole-principal-part}
Pulling the standard simple-pole coordinate on \(\code{OnePoint}\,\C\) back along
the biholomorphism yields a map with a simple-pole principal part at the chosen
point. This is the explicit-input principal-part frontier: prove the pulled-back
coordinate is meromorphic everywhere, extends continuously to \(\infty\) at
\(P\), has chart-local order \(1\), and diverges in norm near \(P\). Open
obligation (\lstinline{sorry}).
Tracking: \ghissue{234}.
\lean{JacobianChallenge.HolomorphicForms.complexSimplePolePrincipalPart_of_biholomorph_onePoint}
\end{lemma}

\begin{lemma}[Single-pole analytic data from simple-pole sphere data]
\label{lem:single-pole-analytic-data-of-simple-pole-to-sphere}
\uses{def:simple-pole-to-sphere-data-of-complex-principal-part}
A concrete \(\code{SimplePoleToSphereData}\) package at \(P\) gives the
single-pole meromorphic analytic data used by the genus-zero meromorphic route.
\lean{JacobianChallenge.HolomorphicForms.singlePoleAnalyticData_of_simplePoleToSphereData}
\leanok
\end{lemma}
\begin{proof}\leanok
The construction fills the meromorphic map fields directly from
\(\code{SimplePoleToSphereData}\), records pole divisor \([P]\), and packages
the degree-one branched-cover data supplied by the simple-pole order.
\end{proof}

\begin{lemma}[Single-pole analytic data from a biholomorphism (open)]
\label{lem:single-pole-analytic-data-of-biholomorph}
\uses{lem:complex-simple-pole-principal-part-biholomorph, def:simple-pole-to-sphere-data-of-complex-principal-part, lem:single-pole-analytic-data-of-simple-pole-to-sphere}
An explicit biholomorphism \(X\simeq\code{OnePoint}\,\C\) gives the prescribed
single-pole analytic meromorphic data at any point \(P\), once the pulled-back
standard coordinate has the complex simple-pole principal part of
Lemma~\ref{lem:complex-simple-pole-principal-part-biholomorph}.
\lean{JacobianChallenge.HolomorphicForms.singlePoleAnalyticData_of_biholomorph_onePoint}
\end{lemma}

\begin{theorem}[Genus-zero single-pole analytic data (open)]
\label{thm:genus-zero-single-pole-analytic-data}
\uses{lem:exists-biholomorph-analyticgenus-zero, lem:single-pole-analytic-data-of-biholomorph}
A compact connected Riemann surface of analytic genus zero admits the
single-pole meromorphic analytic data needed for the prescribed-pole route at
any chosen point \(P\).
\lean{JacobianChallenge.HolomorphicForms.genusZero_singlePoleMeromorphicAnalyticData_nonempty}
This is downstream assembly through the explicit biholomorphic uniformization
frontier and the principal-part frontier. It is deliberately \emph{not} proved
from \(\code{genusZero\_pointRRSection\_*}\) or the fixed-pole
Riemann--Roch section chain, because that route depends on this single-pole
package and would be circular.
\end{theorem}

\begin{lemma}[Pullback of holomorphic 1-forms via biholomorphism]
\label{lem:holomorphic-one-form-pullback-via-biholo}
A homeomorphism \(e:X\to\code{OnePoint}\,\C\) that is biholomorphic
induces a \(\C\)-linear isomorphism of holomorphic 1-form spaces,
\(e^{*}:H^{0}(\code{OnePoint}\,\C,\Omega^{1})\to H^{0}(X,\Omega^{1})\).
\lean{JacobianChallenge.HolomorphicForms.holomorphicOneForm_linearEquiv_of_biholo_to_OnePointCx}
\leanok
\end{lemma}
\begin{proof}
\leanok
Standard pullback of forms via biholomorphism: \(e^*\omega(v):=\omega(de\cdot v)\)
is a cotangent-space pullback, \(\C\)-linear in \(\omega\), and an
isomorphism because \(de\) is invertible at every point.
\end{proof}

\subsection{Chart-trivialisation API for ContMDiffSection}
\label{sec:chart-trivialisation-API}

\begin{lemma}[Identity-chart local representative is smooth]
\label{lem:section-localRepr-identity-chart-contdiff}
The cotangent-bundle section \(\omega.\code{toFun}\) pulled back through
\(\code{identityChart}.\code{symm}\) and evaluated at \(1:\C\) gives a
\(C^{\infty}\) function on \(\C\).
\lean{JacobianChallenge.HolomorphicForms.ContMDiffSection_localRepr_identityChart_contDiff}
\leanok
\end{lemma}
\begin{proof}
\leanok
A \(\code{ContMDiffSection}\) is by definition smooth in any chart of
the bundle's base. In the identity chart of \(\code{OnePoint}\,\C\), the
local representative is a smooth section of the trivial cotangent
bundle on \(\C\); evaluating on the constant tangent vector \(1\) is
a smooth (in fact \(\C\)-linear) operation. \textbf{Mathlib gap}: a
chart-trivialisation API for \code{ContMDiffSection} on the cotangent
bundle is not in v4.28.0.
\end{proof}

\begin{lemma}[Inversion-chart local representative is continuous at zero]
\label{lem:section-localRepr-inversion-chart-continuous-at-zero}
The cotangent-bundle section pulled back through
\(\code{inversionChart}.\code{symm}\) and evaluated at \(1:\C\) is
continuous at \(w=0\).
\lean{JacobianChallenge.HolomorphicForms.ContMDiffSection_localRepr_inversionChart_continuousAt_zero}
\leanok
\end{lemma}
\begin{proof}
\leanok
Same chart-trivialisation argument as
Lemma~\ref{lem:section-localRepr-identity-chart-contdiff}, applied
in the inversion chart instead.
\end{proof}

\subsection{Riemann-Roch \(K-[P]\) Vanishing (Path B)}
\label{sec:rr-k-minus-p-vanishing}

\begin{lemma}[Differentiability of finite lift]
\label{lem:to-finite-fun-mdifferentiable}
\lean{JacobianChallenge.HolomorphicForms.meromorphicFunctionWithDivisors_toFiniteFun_mdifferentiable}
A global complex-valued lift of a meromorphic map with no poles is holomorphic.
\end{lemma}

\begin{proof}\leanok\end{proof}

\begin{lemma}[Germ-order bridge]
\label{lem:germ-order-bridge}
\lean{JacobianChallenge.HolomorphicForms.meromorphicFunctionWithDivisors_order_eq_orderAt}
\leanok
Relates the Divisor order function to the analytic vanishing order.

This is now a direct projection from the
\code{MeromorphicFunctionWithDivisors} order--\code{orderAt} compatibility
field. Constructor-level compatibility is recorded explicitly in the carrier;
the remaining map-to-carrier wrapper proof is the named open obligation
\code{order\_eq\_orderAt\_migration\_obligation}.
\end{lemma}
\begin{proof}\leanok\end{proof}

\begin{lemma}[Value at positive pole]
\label{lem:to-map-eq-infty-of-pole-divisor-pos}
\uses{lem:germ-order-bridge}
\lean{JacobianChallenge.HolomorphicForms.meromorphicFunctionWithDivisors_toMap_eq_infty_of_poleDivisor_pos}
At a positive-order pole, a meromorphic map evaluates to infinity.
\end{lemma}
\begin{proof}
\leanok
Follows by analytic contradiction: finite evaluation implies non-negative germ order.
\end{proof}

\begin{lemma}[Branched cover data of pole degree]
\label{lem:branched-cover-data-of-pole-degree}
\uses{def:branched-degree}
\lean{JacobianChallenge.HolomorphicForms.meromorphicFunctionWithDivisors_branchedCoverDataOfPoleDegree}
A meromorphic function provides valid branched cover data whose branched degree equals the pole divisor degree.
\end{lemma}

\begin{lemma}[Zero divisor degree equals branched degree]
\label{lem:zero-divisor-degree-eq-branched-degree}
\uses{def:branched-degree}
\lean{JacobianChallenge.HolomorphicForms.meromorphicMapToSphere_zeroDivisor_degree_eq_weightedFiberCard_zero}
The degree of the zero divisor of a meromorphic map equals the branched degree of the map. (parked: deep local degree-theory — zeros-with-multiplicity as the weighted fiber over 0; no current owner)
\end{lemma}

\begin{lemma}[Zeros degree equals poles degree]
\label{lem:zeros-degree-eq-poles-degree}
\uses{lem:to-finite-fun-mdifferentiable, lem:to-map-eq-infty-of-pole-divisor-pos, lem:branched-cover-data-of-pole-degree, lem:zero-divisor-degree-eq-branched-degree}
\lean{JacobianChallenge.HolomorphicForms.meromorphicFunctionWithDivisors_degree_zeros_eq_degree_poles}
\notready
For a valid divisor-compatible meromorphic function on a compact Riemann surface, the degree of its zero divisor equals the degree of its pole divisor (open pending the four frontiers).
\end{lemma}
\begin{proof}
Obtain branched cover data, compute degrees.
\end{proof}

\begin{lemma}[Principal divisor degree is zero]
\label{lem:principal-divisor-degree-is-zero}
\uses{lem:zeros-degree-eq-poles-degree}
\lean{JacobianChallenge.HolomorphicForms.meromorphicFunctionWithDivisors_degree_principal_eq_zero}
\notready
The degree of the principal divisor of any non-zero meromorphic function is zero (open pending the four frontiers).
\end{lemma}
\begin{proof}
Subtract pole degree from zero degree.
\end{proof}

\begin{theorem}[Genus zero Riemann-Roch \(K-[P]\) dimension is zero]
\label{thm:genus-zero-riemann-roch-k-minus-p}
\uses{lem:principal-divisor-degree-is-zero}
\lean{JacobianChallenge.HolomorphicForms.genusZero_riemannRoch_K_minus_point_dim_zero}
The cited Lean declaration is currently open pending the RR dimension input via
jc0's \code{genusZero\_pointRiemannRochGermSpaceDimensionInput} provider. The
real statement awaits the \(K-[P]\) germ-space carrier.

Once formalized, the intended statement is: The Riemann-Roch space \(L(K - [P])\) is trivial in genus zero.
\end{theorem}
